\section{Proof of Pairing Selection}
\label{sec:phil-pairing-selection-proof}

\noindent\emph{Proof of \Cref{conj:phil-contrast-formation}.}
Write \(p_i=\pi_{i,i+1}\), encode the assignment by
\(\sigma_i=+1\) when receiver \(i\) selects \(i+1\), and put
\(\varepsilon=\epsilon_{\mathrm L}\). The state remains
\((\mathbf x,\sigma)\in\R^{12}\times\{\pm1\}^4\), with the shared edge
costs of \cref{eq:rg-formation-cycle-cost}. The multiplicity of edge
\(e\) is
\(m_e(\sigma)=1_{\{\sigma_e=+1\}}+1_{\{\sigma_{e+1}=-1\}}\).
In this proof define
\[
 \Omega=I+R_{\mathrm{ev}},\qquad
 \mathsf A_\sigma=H_\sigma+\mathrm T\Omega,\qquad
 H_\sigma=\nabla^2V_\sigma,\qquad
 \mathsf A_\sigma\succeq\mathrm TI.
\]
Let \(\mathsf w_\sigma=\det(\mathsf A_\sigma/\mathrm T)^{-1/2}\),
\(\mathsf Z(p)=\sum_\sigma\pi_\sigma(p)\mathsf w_\sigma\), and
\[
 \mu_p(\mathbf x,\sigma)
 =\frac{\pi_\sigma(p)}{(2\pi)^6\mathsf Z(p)}
   \exp\left[-\frac{\mathbf x^\top\mathsf A_\sigma\mathbf x}
                    {2\mathrm T}\right],\qquad
 \pi_\sigma(p)=\prod_i p_i^{1_{\{\sigma_i=+1\}}}
                         (1-p_i)^{1_{\{\sigma_i=-1\}}}.
\]
Each held-sheet diffusion preserves its Gaussian factor, and a clock ring
resamples \(\sigma\) from its conditional law under \(\mu_p\). Thus
\(\mu_p\) is invariant for the frozen rows, including boundary rows on their
supported sheets. Set
\begin{equation}
 \bar\beta_i(p)=\mu_p(\sigma_i=+1),\qquad
 \mathfrak g_i(p)=\partial_{p_i}\log\mathsf Z(p),\qquad
 \mathscr V_i(p)=\bar\beta_i(p)-p_i
               =p_i(1-p_i)\mathfrak g_i(p).
\label{eq:phil-selection-reduced-field}
\end{equation}
The last identity follows by differentiating the product prior. Since
\(\mathsf Z\) is a positive multilinear polynomial on the closed cube,
\(\mathfrak g\) and \(\mathscr V\) extend analytically to a neighborhood of
that cube and have bounded derivatives there. The well-posedness argument
for \cref{prop:phil-seeded-contrast-formation} applies at every fixed
\(\mathrm T>0\); it gives deterministic continuously differentiable rows and
\begin{equation}
 \min_i\{p_i(t),1-p_i(t)\}\ge\tfrac12e^{-\varepsilon t},\qquad
 \|p(t)-\tfrac12\|_\infty\le\tfrac12\varepsilon t.
\label{eq:phil-selection-row-bounds}
\end{equation}
All constants below may depend on the fixed temperature. Constants in
bounds on a fixed initial law may also depend on its mean, whereas the
window constant in the conclusion will depend only on the temperature.

\subsection*{Frozen convergence and the initial imprint}

Here \(\operatorname{Ent}_\nu(g)=\int g\log(g/\int g d\nu)d\nu\)
for nonnegative \(g\), and total variation is \(\sup_A|q(A)-\nu(A)|\).
For the time-one frozen kernel, the function
\(\mathsf V(\mathbf x)=|\mathbf x|^2\) satisfies
\begin{equation}
 \E_{\mathbf x}|\mathbf X_t|^2
 \le e^{-2\mathrm Tt}|\mathbf x|^2+12(1-e^{-2\mathrm Tt}),\qquad
 P_1^p\mathsf V\le e^{-2\mathrm T}\mathsf V+12.
\label{eq:phil-selection-moment}
\end{equation}
Indeed, jumps leave \(\mathbf X\) unchanged, and the diffusion generator
applied to its squared norm is at most \(24\mathrm T-2\mathrm T|\mathbf X|^2\).
Choose \(R>24/(1-e^{-2\mathrm T})\). For any frozen \(p\), some assignment
\(\sigma_*(p)\) has \(\pi_{\sigma_*}\ge1/16\). Restrict the rate-one
clock to exactly one ring in \([0,1/2]\) and none afterward. Its probability
is \(e^{-1}/2\). Starting with \(|\mathbf x|\le\sqrt R\), the first held
Gaussian motion is in the ball of radius
\(r'=\sqrt R+\sqrt{24\mathrm T}\) at that ring with probability at least
\(1/2\). On the ball, the probability to redraw \(\sigma_*\) is at least
\(16^{-1}e^{-v(r')/\mathrm T}\), where
\(v(r')=\max_{|\mathbf z|\le r',\sigma}V_\sigma(\mathbf z)<\infty\).
The subsequent held Gaussian transition lasts between \(1/2\) and \(1\).
Its means are bounded and its covariance matrices have common positive
lower and finite upper bounds. Its density consequently has a positive
common lower bound on the unit ball. Hence, for some \(\alpha_{\mathrm T}>0\)
independent of \(p\),
\[
 P_1^p((\mathbf x,\sigma),\cdot)
 \ge\alpha_{\mathrm T}\operatorname{Unif}(B_1)
                         \otimes\delta_{\sigma_*(p)}
 \quad\text{when }|\mathbf x|^2\le R.
\]
These are the drift and minorization hypotheses of
\citet[Assumptions 1--2 and Theorems 1.2--1.3]{HairerMattingly2011}.
The explicit dependence in their Theorem 1.3 makes the resulting constants
uniform in the frozen rows, even though the minorizing measure depends on
\(p\). Thus \(\mu_p\) is the unique invariant law and
\begin{equation}
 |P_t^p\phi(\mathbf x,\sigma)-\mu_p\phi|
 \le\mathrm K_{\mathrm T}e^{-\lambda_{\mathrm T}t}
       (1+|\mathbf x|^2)\|\phi\|_\infty.
\label{eq:phil-selection-harris}
\end{equation}
Intermediate times follow by applying the skeleton estimate to
\(P_r^p\phi\), \(0\le r<1\).

Put \(\mathbf v=(1,-1,1,-1)\),
\(f_{\mathbf v}(\sigma)=\frac12\sum_i\mathbf v_i\sigma_i\), and define
\begin{equation}
 \mathsf I_i(\mathsf m_0)
 =\int_0^\infty[q_t^0(\sigma_i=+1)-\tfrac12]dt,\qquad
 \mathsf j_{\mathrm T}(\mathsf m_0)
 =\mathbf v\cdot\mathsf I(\mathsf m_0).
\label{eq:phil-selection-imprint}
\end{equation}
At uniform rows, the stationary law is invariant under cycle rotations and
reflections, so each stationary row is \(1/2\) and
\(\mu_{1/2}f_{\mathbf v}=0\). Estimate
\cref{eq:phil-selection-harris} proves absolute convergence of these
integrals. It also gives a Poisson integral
\begin{equation}
 \mathfrak R_{\mathrm T}(\mathbf x,\sigma)
 =\int_0^\infty P_t^{1/2}f_{\mathbf v}(\mathbf x,\sigma)dt,\qquad
 |\mathfrak R_{\mathrm T}(\mathbf x,\sigma)|
 \le\mathrm K_{\mathrm T}(1+|\mathbf x|^2).
\label{eq:phil-selection-poisson}
\end{equation}
Convolution of this integral with the initial Gaussian expresses
\(\mathsf j_{\mathrm T}\) as an entire function of its complexified mean:
for every compact set of complex means, all differentiated integrands are
bounded by an integrable polynomial times a Gaussian. The same holds for
\(\mathsf I\). A rotation by one changes the sign of
\(\mathsf j_{\mathrm T}\).

For \(t\le1/(10\varepsilon)\), all row entries are in
\([\mathrm c_0,1-\mathrm c_0]\), \(\mathrm c_0=0.45\). A posterior receiver probability has derivative
at most \(1/[4\mathrm c_0(1-\mathrm c_0)]\le1/(2\mathrm c_0)\) with respect to its prior entry.
Coupling the common clocks and their four posterior rows therefore gives
\[
 \|q_t^\varepsilon-q_t^0\|_{\mathrm{TV}}
 \le\frac2{\mathrm c_0}\int_0^t\|p(s)-\tfrac12\|_\infty ds
 \le\frac{\varepsilon t^2}{2\mathrm c_0}.
\]
Variation of constants in the learning equation, followed by
\cref{eq:phil-selection-harris}, gives
\begin{equation}
 \left|p(t)-\tfrac12-\varepsilon\mathsf I(\mathsf m_0)\right|_\infty
 \le\mathrm K_{\mathrm T,\mathsf m_0}
       (\varepsilon e^{-\lambda_{\mathrm T}t}+\varepsilon^2t)
       +\frac{\varepsilon^2t^3}{6\mathrm c_0}.
\label{eq:phil-selection-imprint-expansion}
\end{equation}
In particular, at \(t_1=\varepsilon^{-1/4}\),
\(p(t_1)-1/2=\varepsilon\mathsf I(\mathsf m_0)+o(\varepsilon)\).

\subsection*{Nonvanishing of the matching imprint}

Take \(\mathbf u_\circ=(0_4,0_4,(0,1,0,-1))\), whose two residual vectors
are \((0,-1,0,1)\). In the noiseless switched motion the subspace
\(\bar k=\bar m\) is invariant. After removing the scalar factor
\(e^{-\mathrm Tt}\), its common residual vector satisfies
\begin{equation}
 \dot\rho=-\mathsf G M_\sigma\rho,\qquad
 \mathsf G=4I-A_{C_4},\qquad
 M_\sigma=\operatorname{diag}(m_e(\sigma)).
\label{eq:phil-selection-tied-residual}
\end{equation}
Here the cyclic difference matrix contributes
\(2I-A_{C_4}\) and the shared-link drift contributes \(2I\).
Of the sixteen assignments, six have \(m_1=m_3\). Their first held motion
is \(e^{-ct}(0,-1,0,1)\), with \(c=4m_1\in\{0,4,8\}\).
At their first clock ring every receiver prefers the even-edge matching
\(S^{(0)}\). Holding that matching then leaves the odd residuals constant
and the even residuals zero in the scalar-free coordinates.

Four further assignments have multiplicity patterns
\[
 (1,1,2,0),\quad(1,0,2,1),\quad(2,0,1,1),\quad(2,1,1,0).
\]
The signed swap of coordinates \(1,3\) and the rotation by two preserve
the even-edge matching and carry the initial residual to itself or its
negative. They reduce these four cases to \((1,1,2,0)\). Write its
scalar-free coordinates as \(\rho=(-a,-b,-d,h)\). Then
\[
 \begin{pmatrix}\dot a\\\dot b\\\dot d\end{pmatrix}
 =\begin{pmatrix}-4&1&0\\1&-4&2\\0&1&-8\end{pmatrix}
   \begin{pmatrix}a\\b\\d\end{pmatrix},\qquad
 (a,b,d)(0)=(0,1,0),\qquad \dot h=-(a+2d),\ h(0)=1.
\]
The nonnegative cone is invariant, and the negative row sums give
\(\max(a,b,d)\le1\), hence \(a\le1/4\), \(d\le1/8\).
The half-spaces
\(a-3b/4-11d/25\le0\) and \(d-b/4\le0\) are invariant: their
outward derivatives on the corresponding boundary are
\(-b/400-7d/100\) and \(-a/4-b/8\), respectively. Consequently
\(a\le43b/50\), \(d\le b/4\), and \(b\ge e^{-4t}\).
The attended block is similar to a negative definite symmetric matrix
because \(\mathsf G\succ0\). Solving
\(\rho(\infty)=\rho(0)-\mathsf G z\), with \(z_3=0\) and the first
three final coordinates zero, gives \(h(\infty)=6/7\). Thus
\(h\ge6/7\).

If the first ring occurs at \(t_{\mathrm r}\) and redraws \(S^{(0)}\),
the subsequent scalar-free held motion, in elapsed time \(r\), is
\[
 \begin{aligned}
 a_r&=a_0e^{-8r},&d_r&=d_0e^{-8r},\\
 b_r&=b_0+\tfrac14(a_0+d_0)(1-e^{-8r}),&
 h_r&=h_0-\tfrac14(a_0+d_0)(1-e^{-8r})\ge\tfrac{171}{224}.
 \end{aligned}
\]
Every adjacent odd-edge absolute residual exceeds its even-edge neighbor
by at least \((7/50)e^{-4t_{\mathrm r}}\). Combining these four cases
with the preceding six, the deterministic motion \(\mathbf Y\) from
\(s\mathbf u_\circ\) has, throughout the hypothesized matching hold,
\begin{equation}
 \sqrt{\mathcal D_{\mathrm{odd}}(\mathbf Y_t)}
 -\sqrt{\mathcal D_{\mathrm{even}}(\mathbf Y_t)}
 \ge\frac7{50}s e^{-\mathrm Tt-8t_{\mathrm r}},
 \qquad t\ge t_{\mathrm r}.
\label{eq:phil-selection-lock-margin}
\end{equation}
There are therefore ten good initial assignments. For the remaining six
only \(f_{\mathbf v}\ge-2\) will be used.

Set \(h_s=(\log s)/\mathrm T\), for \(s>e\), and write the initial
configuration as \(s\mathbf u_\circ+\xi\), \(\xi\sim\mathcal N(0,I)\).
Along the actual assignment path decompose
\(\mathbf X=\mathbf Y+\mathbf E+\mathbf Z\), where the first two terms
solve its switched deterministic equation from \(s\mathbf u_\circ\) and
\(\xi\), respectively, while
\[
 d\mathbf Z=-\mathsf A_{\sigma_t}\mathbf Zdt+\sqrt{2\mathrm T}dW_t,
 \qquad\mathbf Z_0=0.
\]
All three are adapted differential equations, and
\(|\mathbf E_t|\le e^{-\mathrm Tt}|\xi|\).
It\^o's formula yields
\[
 \mathscr L e^{|\mathbf Z|^2/4}
 \le\mathrm T(6-|\mathbf Z|^2/4)e^{|\mathbf Z|^2/4}
 \le\mathrm Te^5.
\]
Stopping at radius \(\log s\), and using
\(\E e^{|\xi|^2/4}=64\), bounds the probability that either
\(|\xi|>\log s\) or \(\sup_{t\le h_s}|\mathbf Z_t|>\log s\) by
\((65+\mathrm Te^5h_s)e^{-(\log s)^2/4}\).
On its complement each square-root edge cost changes from its
\(\mathbf Y\) value by at most \(2\mathrm C_N\log s\), where
\(\mathrm C_N\) is the largest norm of a residual-pair map divided by
\(\sqrt2\).

Define
\[
 L_s(t_{\mathrm r})=h_s-\frac8{\mathrm T}t_{\mathrm r}
                         -\frac2{\mathrm T}\log\log s,\qquad
 B_s=(L_s(t_{\mathrm r})-t_{\mathrm r})_+.
\]
For a good initial assignment, \cref{eq:phil-selection-lock-margin}
gives a square-root gap at least \((7/50)(\log s)^2\) on
\([t_{\mathrm r},L_s]\). For sufficiently large \(s\), while both noise bounds hold,
the actual energy gap is at least
\(\kappa_0(\log s)^4\), \(\kappa_0=(7/50)^2/4\).
Each receiver's wrong choice at a ring then has probability at most
\(e^{-\kappa_0(\log s)^4/\mathrm T}\). Apply this bound at the first
ring, then use the rate-one compensator up to the first wrong redraw or
exit from the noise bound, stopping at \(L_s\wedge h_s\).
The endpoint \(L_s\) is known after the first ring, and the gap estimate
holds until that first failure. Thus the matching holds throughout the
nonempty interval \([t_{\mathrm r},L_s]\) except on an event of
probability at most
\begin{equation}
 \delta_s=(65+\mathrm Te^5h_s)e^{-(\log s)^2/4}
       +4(1+h_s)e^{-\kappa_0(\log s)^4/\mathrm T},
 \qquad h_s\delta_s\longrightarrow0.
\label{eq:phil-selection-lock-failure}
\end{equation}
Since \(\E t_{\mathrm r}=1\), a good initial assignment contributes at
least
\[
 -2h_s+4\E B_s-4h_s\delta_s
 \ge2h_s-4(1+8/\mathrm T)-(8/\mathrm T)\log\log s-4h_s\delta_s
\]
to \(\E\int_0^{h_s}f_{\mathbf v}(\sigma_t)dt\). The other six contribute
at least \(-2h_s\). Averaging the sixteen uniform initial assignments gives
\[
 \E\int_0^{h_s}f_{\mathbf v}(\sigma_t)dt
 \ge\tfrac12h_s-\frac5{\mathrm T}\log\log s-O_{\mathrm T}(1).
\]
By \cref{eq:phil-selection-moment},
\(\E|\mathbf X_{h_s}|^2\le14\). The Markov property and
\cref{eq:phil-selection-poisson} bound the integral after \(h_s\) in
absolute value by a constant independent of \(s\). Therefore
\begin{equation}
 \mathsf j_{\mathrm T}(s\mathbf u_\circ)
 \ge\frac{\log s}{2\mathrm T}
       -\frac5{\mathrm T}\log\log s-O_{\mathrm T}(1).
\label{eq:phil-selection-nonvanishing}
\end{equation}
This proves that \(\mathsf j_{\mathrm T}\) is not identically zero at any
fixed \(\mathrm T>0\). Its zero set is Lebesgue-null. For completeness,
this last assertion follows by induction on dimension: a nonzero analytic
function in one variable has isolated zeros; in several variables at least
one coefficient \(\partial_n^kf(\mathbf x',0)\) is nonzero as an analytic
function, or analyticity would make \(f\) vanish everywhere. Outside the
null zero set of that coefficient, each one-dimensional slice has isolated
zeros, and Fubini's theorem concludes the induction.

\subsection*{Entropy balance for the evolving recognition law}

Fix a finite horizon \(h\). All rows are bounded away from zero on this
interval by \cref{eq:phil-selection-row-bounds}. Write \(q_\sigma(t,\mathbf x)\)
for the evolving density, \(Q=\sum_\sigma q_\sigma\), and
\(b_\sigma(t,\mathbf x)=\mu_{p(t)}(\mathbf x,\sigma)/
 \sum_\eta\mu_{p(t)}(\mathbf x,\eta)\).
For the forward Ornstein--Uhlenbeck semigroup \(K_\sigma\), the last-proposal
formula is
\begin{equation}
 q_\sigma(t)=e^{-t}K_\sigma(t)q_\sigma(0)
   +\int_0^te^{-(t-r)}K_\sigma(t-r)[b_\sigma(r)Q(r)]dr.
\label{eq:phil-selection-forward}
\end{equation}
Iterating this positive formula, and bounding each proposal factor by one,
reduces each term to a Gaussian motion along a deterministic sequence of
sheets. Put \(M=\max_\sigma\|\mathsf A_\sigma\|\).
Its means have norm at most \(|\mathsf m_0|\), and its covariances lie between
\(e^{-2Mh}I\) and \((1+2\mathrm Th)I\). Summing the at most \(16^n\)
sequences on each time simplex, of volume \(t^n/n!\), gives a finite upper
Gaussian envelope. The no-proposal term gives a lower Gaussian envelope.
Consequently, uniformly on \([0,h]\),
\begin{equation}
 0<\mathrm c_he^{-\mathrm C_h(1+|\mathbf x|^2)}
 \le q_\sigma(t,\mathbf x)\le\mathrm C_he^{-\mathrm c_h|\mathbf x|^2}.
\label{eq:phil-selection-density-envelopes}
\end{equation}
This bound does not assert Gaussianity conditional on a state-dependent
assignment history.

Differentiating one Ornstein--Uhlenbeck kernel against a function bounded
by a polynomial times a Gaussian gives a bound of the same type with
factor \((t-r)^{-1/2}\). To check it, use the Gaussian kernel covariance
bounds proportional to \(t-r\), absorb its linear derivative into a
Gaussian with twice the covariance, and complete the square in the
convolution. This singularity is integrable in
\cref{eq:phil-selection-forward}, proving a first-derivative Gaussian
bound. Moreover
\[
 \nabla b_\sigma=-\frac{b_\sigma}{\mathrm T}
       \left(H_\sigma\mathbf x-\sum_\eta b_\eta H_\eta\mathbf x\right),
 \qquad |\nabla b_\sigma|\le\mathrm K_{\mathrm T}|\mathbf x|.
\]
The source \(b_\sigma Q\) therefore has a first derivative with the same
type of bound. For the second derivative, transfer one derivative to this
source using
\(\nabla K_\sigma(r)g=e^{r\mathsf A_\sigma}K_\sigma(r)\nabla g\),
and differentiate the other kernel factor. The singularity is again only
\(r^{-1/2}\). The explicit initial Gaussian term supplies the bounds at
zero. The forward equation obtained by taking a time difference quotient
in \cref{eq:phil-selection-forward} then bounds the time derivative too:
\begin{equation}
 |\nabla q_\sigma|+|\nabla^2q_\sigma|+|\partial_tq_\sigma|
 \le\mathrm C_h(1+|\mathbf x|)^k e^{-\mathrm c_h|\mathbf x|^2}.
\label{eq:phil-selection-density-derivatives}
\end{equation}
The derivative transfer follows directly by integration by parts in the
Gaussian transition density; its inverse mean map is uniformly bounded
on this finite interval. These estimates justify the differentiations,
without assuming finite Fisher information.

Set \(f_\sigma=q_\sigma/\mu_{p(t),\sigma}\) and
\(\mathsf H(t)=\KL(q_t\Vert\mu_{p(t)})\).
The logarithms of both densities grow at most quadratically by
\cref{eq:phil-selection-density-envelopes}, and
\(\partial_t\log\mu_p\) is bounded in \(\mathbf x\), since only its prior
weight and normalizer vary. Together with
\cref{eq:phil-selection-density-derivatives}, this proves that
\(\mathsf H\) is absolutely continuous on \([0,h]\) and permits its
integral to be differentiated.
To integrate the diffusion by parts, first multiply by a smooth radial
cutoff \(\chi_R\), equal to one on \(B_R\), zero outside \(B_{2R}\),
with \(|\nabla\chi_R|\le\mathrm C/R\), increasing to one.
On each compact cylinder the densities have positive lower bounds, so
\(\log f_\sigma\) is a legitimate test function. The boundary term is
bounded by
\[
 \frac{\mathrm C_h}{R}\int_{R\le|\mathbf x|\le2R}
       (1+|\mathbf x|)^{k+3}e^{-\mathrm c_h|\mathbf x|^2}d\mathbf x
 \longrightarrow0.
\]
Detailed balance \(\mu_\sigma b_\eta=\mu_\eta b_\sigma\) symmetrizes the
finite jump sum. Both resulting production integrands are nonnegative,
so monotone convergence removes the cutoff and simultaneously proves
finiteness and local time integrability of the production. The result is
\begin{equation}
 \begin{aligned}
 \mathsf H'(t)&=-\mathscr D_{p(t)}(f)
               -\sum_\sigma\int q_\sigma\partial_t\log\mu_{p(t),\sigma},\\
 \mathscr D_p(f)&=\mathrm T\sum_\sigma\int\mu_{p,\sigma}
                         \frac{|\nabla f_\sigma|^2}{f_\sigma}\\
 &\quad+\frac12\sum_{\sigma,\eta}\int\mu_{p,\sigma}b_\eta
                  (f_\eta-f_\sigma)(\log f_\eta-\log f_\sigma).
 \end{aligned}
\label{eq:phil-selection-entropy-identity}
\end{equation}
The equality holds almost everywhere and in integrated form on every finite
interval. The jump term is integrable already from the Gaussian envelopes
and quadratic logarithmic bounds; thus the cutoff argument does not
presuppose a bound on the diffusion production. This also proves
\(\mathsf H(0)<\infty\).

\subsection*{Entropy contraction at frozen rows}

In the original twelve-dimensional coordinates put
\(\mathsf B_\sigma=\mathsf A_\sigma/\mathrm T\), the Gaussian sheet
precision, and set
\[
 \kappa=\min_\sigma\lambda_{\min}(\mathsf B_\sigma)>0,\qquad
 M=\max_\sigma\|\mathsf B_\sigma\|,\qquad
 q_\sigma^*=\mu_p(\sigma)
 =\frac{\pi_\sigma\mathsf w_\sigma}{\mathsf Z},
 \qquad q_{\min}^*=\min_\sigma q_\sigma^*.
\]
Let \(\gamma_\sigma=\mathcal N(0,\mathsf B_\sigma^{-1})\), and, for a
real function \(u\), define
\[
 \mathscr E_c(u)=\mathrm T\sum_\sigma\int\mu_{p,\sigma}|\nabla u_\sigma|^2,
 \qquad
 \mathscr E_j(u)=\frac12\sum_{\sigma,\eta}\int\mu_{p,\sigma}b_\eta
                                    (u_\sigma-u_\eta)^2.
\]
The Gaussian inequalities needed below follow from the explicit
Ornstein--Uhlenbeck semigroup with invariant law \(\gamma_\sigma\).
Its gradient commutes as
\(\nabla P_tg=e^{-t\mathsf B_\sigma}P_t\nabla g\). Differentiating
entropy and applying Cauchy--Schwarz inside this kernel gives
\[
 -\frac d{dt}\operatorname{Ent}_{\gamma_\sigma}(P_tg)
 =\int\frac{|\nabla P_tg|^2}{P_tg}d\gamma_\sigma
 \le e^{-2\kappa t}\int\frac{|\nabla g|^2}{g}d\gamma_\sigma.
\]
For bounded positive smooth \(g\), the explicit Gaussian formula and
dominated convergence make the entropy tend to zero. Integration, with
\(g=u^2\), gives
\begin{equation}
 \operatorname{Ent}_{\gamma_\sigma}(u^2)
 \le\frac2\kappa\int|\nabla u|^2d\gamma_\sigma,
 \qquad
 \operatorname{Var}_{\gamma_\sigma}(u)
 \le\frac1\kappa\int|\nabla u|^2d\gamma_\sigma.
\label{eq:phil-selection-gaussian-inequalities}
\end{equation}
The second estimate follows by differentiating variance in the same
semigroup. Positive regularization, truncation, and smooth cutoffs extend
both estimates to their form domains: the cutoff errors tend to zero by
integrability of the squared function and its gradient, and the entropy
passes to the limit by its variational formula.

Replace \(u\) by \(|u|\), which decreases both energies, and normalize
\(\int u^2d\mu_p=1\). Put
\(z_\sigma=(\int u_\sigma^2d\gamma_\sigma)^{1/2}\).
The entropy chain rule gives
\[
 \operatorname{Ent}_{\mu_p}(u^2)
 =\sum_\sigma q_\sigma^*
             \operatorname{Ent}_{\gamma_\sigma}(u_\sigma^2)
       +\operatorname{Ent}_{q^*}(z^2).
\]
Here \(0\le z_\sigma\le R=(q_{\min}^*)^{-1/2}\).
The function \(\psi(z)=z^2\log z^2-z^2+1\) satisfies
\(\psi(1)=\psi'(1)=0\) and \(\psi''(z)=4+4\log z\), whence
\(\psi(z)\le2(1+\log R)(z-1)^2\) on \([0,R]\).
Since \(\sum q_\sigma^*z_\sigma^2=1\),
\[
 \operatorname{Ent}_{q^*}(z^2)
 \le L_p\operatorname{Var}_{q^*}(z),\qquad
 L_p=4+2\log(1/q_{\min}^*).
\]
Indeed \(\sum q_\sigma^*(z_\sigma-1)^2
=2(1-\E_{q^*}z)\le2\operatorname{Var}_{q^*}(z)\).

Fix a ball \(B=B_\rho\) in these coordinates, and define
\[
 m_\sigma=|B|^{-1}\int_Bu_\sigma,
 \quad m_* =\sum_\sigma\pi_\sigma m_\sigma,
 \quad X_\sigma=
 \left[\frac{(2\pi)^6\mathsf w_\sigma
                   e^{\rho^2\|\mathsf B_\sigma\|/2}}{|B|}-1\right]_+,
 \quad B_{\mathrm T}=\max_\sigma(1+\sqrt{X_\sigma})^2.
\]
Variance about any center is an upper bound on the variance, so
\[
 \operatorname{Var}_{q^*}(z)
 \le2\sum_\sigma q_\sigma^*(z_\sigma-m_\sigma)^2
     +2\sum_\sigma q_\sigma^*(m_\sigma-m_*)^2.
\]
The first sum is at most \((B_{\mathrm T}/(\mathrm T\kappa))\mathscr E_c(u)\):
\(|z_\sigma-\gamma_\sigma u_\sigma|\le
\sqrt{\operatorname{Var}_{\gamma_\sigma}(u_\sigma)}\), and
Cauchy--Schwarz bounds \(|m_\sigma-\gamma_\sigma u_\sigma|\) by the same
standard deviation times
\(\sqrt{\chi^2(\operatorname{Unif}(B)\Vert\gamma_\sigma)}
\le\sqrt{X_\sigma}\).
For the second sum, Jensen's inequality gives
\[
 \sum_\sigma q_\sigma^*(m_\sigma-m_*)^2
 \le\frac{\mathsf w_{\max}}
              {\mathsf Z|B|}
       \sum_{\sigma,\eta}\pi_\sigma\pi_\eta
                          \int_B(u_\sigma-u_\eta)^2.
\]
On \(B\), the explicit Gaussian densities give
\[
 \mu_{p,\sigma}b_\eta
 \ge\frac{\pi_\sigma\pi_\eta e^{-\rho^2M}}
              {(2\pi)^6\mathsf Z}.
\]
The normalizer therefore cancels, and the preceding sum is at most
\(2(2\pi)^6\mathsf w_{\max}e^{\rho^2M}
 \mathscr E_j(u)/|B|\).
Combining these bounds proves
\[
 \operatorname{Ent}_{\mu_p}(u^2)
 \le\frac2{\mathrm T\kappa}(1+L_pB_{\mathrm T})\mathscr E_c(u)
    +\frac{4L_p(2\pi)^6\mathsf w_{\max}e^{\rho^2M}}{|B|}
                                                       \mathscr E_j(u).
\]
For \(u=\sqrt f\), the production is at least
\(4(\mathscr E_c(u)+\mathscr E_j(u))\). Its jump part uses
\((a-b)(\log a-\log b)\ge4(\sqrt a-\sqrt b)^2\), which follows by
Cauchy--Schwarz on the integral of \(r^{-1/2}\) between \(a\) and \(b\).
The inequality extends to the form domain by the preceding approximations.
Finally,
\[
 q_{\min}^*\ge p_{\min}^4
                 \frac{\mathsf w_{\min}}
                      {\mathsf w_{\max}},\qquad
 p_{\min}=\min_i\{p_i,1-p_i\}.
\]
Thus there are finite positive constants depending only on \(\mathrm T\)
with
\begin{equation}
 \mathscr D_p(f)\ge\lambda(\mathrm T,p)
                         \operatorname{Ent}_{\mu_p}(f),\qquad
 \lambda(\mathrm T,p)\ge
 \frac{\mathrm c_{\mathrm T}}
      {\mathrm K_{\mathrm T}+8\log(1/p_{\min})}.
\label{eq:phil-selection-entropy-contraction}
\end{equation}
The factor \(\mathrm T\) in \(\mathscr E_c\) is the same diffusion
coefficient as in \cref{eq:phil-selection-entropy-identity}.

\subsection*{The matching branch of the small-temperature landscape}

The needed sheet weights can be computed using two four-dimensional
determinants. Let \(D\) be the cyclic difference matrix, and
\(C_\sigma=\operatorname{diag}(m_e(\sigma))\). In the orthogonal
coordinates \((\bar k+\bar m)/\sqrt2\),
\((\bar k-\bar m)/\sqrt2\), and \(\omega\), the energy splits into
\(\frac12(Du-\sqrt2\omega)^\top C_\sigma(Du-\sqrt2\omega)
 +\frac12(Dv)^\top C_\sigma Dv\), while the discrepancy coordinate has
reference precision \(5I\). The determinant identity for the first block
then gives
\begin{equation}
 \det\mathsf A_\sigma
 =\mathrm T^4
  \det[\mathrm TI+C_\sigma(DD^\top+2I)]
  \det[5\mathrm TI+D^\top C_\sigma D].
\label{eq:phil-selection-sheet-determinants}
\end{equation}
The two matching assignments, eight end-path assignments, four middle-path
assignments, and two oriented cycles have the respective determinants
\[
 \begin{aligned}
 \mathsf d_M&=25\mathrm T^8(\mathrm T+8)^2(5\mathrm T+4)^2,\\
 \mathsf d_E&=5\mathrm T^6
   (\mathrm T^3+16\mathrm T^2+77\mathrm T+112)
   (125\mathrm T^3+200\mathrm T^2+85\mathrm T+8),\\
 \mathsf d_N&=5\mathrm T^6(\mathrm T+4)(5\mathrm T+2)
   (\mathrm T^2+12\mathrm T+28)(25\mathrm T^2+30\mathrm T+4),\\
 \mathsf d_C&=5\mathrm T^5(\mathrm T+2)(\mathrm T+4)^2(\mathrm T+6)
                                     (5\mathrm T+2)^2(5\mathrm T+4).
 \end{aligned}
\]
These follow in \cref{eq:phil-selection-sheet-determinants} from the
multiplicity representatives \((2,0,2,0)\), \((2,1,1,0)\),
\((1,2,1,0)\), and \((1,1,1,1)\), respectively.
Normalize the matching weight to one and write
\[
 \mathsf r_E=\sqrt{\mathsf d_M/\mathsf d_E},\qquad
 \mathsf r_N=\sqrt{\mathsf d_M/\mathsf d_N},\qquad
 \mathsf r_C=\sqrt{\mathsf d_M/\mathsf d_C}.
\]
Their leading terms are
\(\mathsf r_E,\mathsf r_N=\sqrt{40/7}\mathrm T+O(\mathrm T^2)\) and
\(\mathsf r_C=\sqrt{5/3}\mathrm T^{3/2}(1+O(\mathrm T))\).
Choose \(\mathrm T_0>0\) so that \(0<\mathsf r_E,\mathsf r_N,\mathsf r_C\le1/8\) for
\(0<\mathrm T\le\mathrm T_0\).

Use matching coordinates
\(x=(p_0,1-p_1,p_2,1-p_3)\). The normalized evidence is the
multilinear polynomial whose vertex weights are one at \(\mathbf0,\mathbf1\),
\(\mathsf r_E\) for subsets of cardinality one or three, \(\mathsf r_N\) for the four adjacent
pairs, and \(\mathsf r_C\) for the two opposite pairs. Complementation preserves
these weights, so the centered reduced field is odd. With
\(D_0=1+4\mathsf r_E+2\mathsf r_N+\mathsf r_C\), its Jacobian at \(x=\mathbf1/2\) is
\[
 \frac1{D_0}
 \begin{pmatrix}
 0&1-\mathsf r_C&1-2\mathsf r_N+\mathsf r_C&1-\mathsf r_C\\
 1-\mathsf r_C&0&1-\mathsf r_C&1-2\mathsf r_N+\mathsf r_C\\
 1-2\mathsf r_N+\mathsf r_C&1-\mathsf r_C&0&1-\mathsf r_C\\
 1-\mathsf r_C&1-2\mathsf r_N+\mathsf r_C&1-\mathsf r_C&0
 \end{pmatrix}.
\]
For example, the adjacent Hessian entry of the evidence is
\((1-\mathsf r_C)/2\), its uniform value is \(D_0/8\), and the row mobility is
\(1/4\), giving the displayed entry. The eigenvalues are
\begin{equation}
 \alpha=\frac{3-2\mathsf r_N-\mathsf r_C}{D_0}>0,\qquad
 \beta=\frac{-1-2\mathsf r_N+3\mathsf r_C}{D_0}<0,\qquad
 \eta=\frac{-1+2\mathsf r_N-\mathsf r_C}{D_0}<0\quad\text{twice}.
\label{eq:phil-selection-saddle-spectrum}
\end{equation}
The positive eigendirection is \(\mathbf1\) in matching coordinates,
or \(\mathbf v\) in the original rows. In fact the chosen bounds give
\(\alpha\ge7/5\), \(\beta\le-1/3\), and \(\eta\le-2/5\).

Cyclic invariance makes the matching diagonal \(x=z\mathbf1\) invariant.
Writing \(w=z(1-z)\), its normalized evidence is
\[
 Z_{\mathrm{diag}}(z)
 =1+4(\mathsf r_E-1)w+(2-8\mathsf r_E+4\mathsf r_N+2\mathsf r_C)w^2,
 \qquad
 \dot z=\frac{z(1-z)}{4Z_{\mathrm{diag}}(z)}Z_{\mathrm{diag}}'(z).
\]
The derivative is \((1-2z)[4(\mathsf r_E-1)+2(2-8\mathsf r_E+4\mathsf r_N+2\mathsf r_C)w]\).
The bracket is negative at both endpoints \(w=0,1/4\), where its values
are \(4(\mathsf r_E-1)\) and \(-3+2\mathsf r_N+\mathsf r_C\). It is affine in \(w\), hence negative
throughout. Thus the diagonal motion tends to \(\mathbf1\) from
\(z>1/2\), and to \(\mathbf0\) from \(z<1/2\), reaching every prescribed
endpoint neighborhood in finite slow time. At either vertex the Jacobian
is \((\mathsf r_E-1)I\): a single receiver flip changes weight one to \(\mathsf r_E\), the
mobility vanishes, and its diagonal derivative supplies the factor
\(-(1-\mathsf r_E)\). Both vertices are therefore sinks with positive inward rate
\(\nu=1-\mathsf r_E\).

\subsection*{Tracking, escape, and selection}

Let \(a_i(t)=q_t(\sigma_i=+1)\) and
\(\Delta_i(t)=a_i(t)-\bar\beta_i(p(t))\).
Differentiating the normalized product prior gives
\[
 \partial_{p_i}\log\mu_p
 =\frac{1_{\{\sigma_i=+1\}}-\bar\beta_i(p)}{p_i(1-p_i)}.
\]
Inserting the learning equation and
\cref{eq:phil-selection-reduced-field} in
\cref{eq:phil-selection-entropy-identity} yields
\begin{equation}
 \mathsf H'
 =-\mathscr D_p(f)-\varepsilon\sum_i\mathfrak g_i(p)\Delta_i
                  -\varepsilon\sum_i\frac{\Delta_i^2}{p_i(1-p_i)}
 \le-\lambda(\mathrm T,p)\mathsf H
      +\varepsilon\|\mathfrak g(p)\|_1\sqrt{\mathsf H/2}.
\label{eq:phil-selection-learning-entropy}
\end{equation}
The last bound uses \(|\Delta_i|\le\operatorname{TV}(q_t,\mu_p)
\le\sqrt{\mathsf H/2}\). The nonpositive square is kept until this step;
it removes any inverse-row factor from the forcing.
Writing \(\Lambda(t)=\int_0^t\lambda(\mathrm T,p(r))dr\), comparison
for \(\sqrt{\mathsf H+\delta^2}\), followed by \(\delta\downarrow0\), gives
\begin{equation}
 \sqrt{\mathsf H(t)}
 \le e^{-\Lambda(t)/2}\sqrt{\mathsf H(0)}
    +\frac{\varepsilon}{2\sqrt2}
       \int_0^te^{-(\Lambda(t)-\Lambda(r))/2}
                         \|\mathfrak g(p(r))\|_1dr.
\label{eq:phil-selection-tracking}
\end{equation}
The bounded fitness, \cref{eq:phil-selection-row-bounds}, and
\cref{eq:phil-selection-entropy-contraction} show that, on
\(t\le2\mathrm C_{\mathrm T}\varepsilon^{-1}\log(1/\varepsilon)\),
\begin{equation}
 \operatorname{TV}(q_t,\mu_{p(t)})
 \le e^{-\Lambda(t)/2}\sqrt{\mathsf H(0)/2}
       +O_{\mathrm T}(\varepsilon\log(1/\varepsilon)).
\label{eq:phil-selection-global-tracking}
\end{equation}
Indeed \(\lambda\) is bounded below on this horizon by a positive
constant divided by \(1+\log(1/\varepsilon)\).

For the saddle a finer bound is needed. Let
\(\|\mathfrak g(p)-\mathfrak g(p')\|_1
 \le L_{\mathrm T}\|p-p'\|_\infty\), and suppose
\(\lambda\ge\lambda_\flat>0\) on \([s_0,t]\).
The estimate \(|\dot p_i|\le\varepsilon\) in
\cref{eq:phil-selection-tracking} gives
\begin{equation}
 \operatorname{TV}(q_t,\mu_{p(t)})
 \le e^{-\lambda_\flat(t-s_0)/2}\sqrt{\mathsf H(s_0)/2}
   +\frac{\varepsilon}{2\lambda_\flat}\|\mathfrak g(p(t))\|_1
   +\frac{L_{\mathrm T}\varepsilon^2}{\lambda_\flat^2}.
\label{eq:phil-selection-refined-tracking}
\end{equation}
The constants follow by integrating \(1\) and \(t-r\) against
\(e^{-\lambda_\flat(t-r)/2}\).

Fix \(0<\mathrm T\le\mathrm T_0\) and a mean with
\(\mathsf j_{\mathrm T}(\mathsf m_0)\ne0\), and put
\(\mathbf e=\mathbf v/2\). In slow time
\(\mathfrak s=\varepsilon t\), write
\(y=p-1/2=a\mathbf e+z\), \(z\perp\mathbf e\).
By \cref{eq:phil-selection-imprint-expansion}, at
\(t_1=\varepsilon^{-1/4}\),
\[
 a(t_1)=\varepsilon(\mathsf j_{\mathrm T}(\mathsf m_0)/2+o(1)),
 \qquad |z(t_1)|=O(\varepsilon).
\]
On \([0,t_1]\), symmetry gives \(\mathfrak g(1/2)=0\), so
\(\|\mathfrak g(p(t))\|_1\le L_{\mathrm T}\varepsilon t/2\).
The bound in \cref{eq:phil-selection-tracking} gives
\(\sqrt{\mathsf H(t_1)}=O(\varepsilon^{7/4})\).
Consequently, until exit from any sufficiently small fixed ball,
\cref{eq:phil-selection-refined-tracking} implies
\begin{equation}
 |\Delta|_2\le\kappa_\varepsilon|y|+b_\varepsilon,
 \qquad \kappa_\varepsilon=O(\varepsilon),\quad
 b_\varepsilon=O(\varepsilon^{7/4})=o(|a(t_1)|).
\label{eq:phil-selection-relative-error}
\end{equation}
The constants here may depend on the fixed mean.

Let \(\beta_\perp=\max(\beta,\eta)<0\) and
\(\gamma=\alpha-\beta_\perp>0\).
Oddness and diagonal invariance of the analytic reduced field imply,
for \(|y|\le r\),
\[
 |\langle\mathscr V(1/2+y),\mathbf e\rangle-\alpha a|
       \le\mathrm C|y|^3,\qquad
 |P_\perp\mathscr V(1/2+y)-D\mathscr V(1/2)z|
       \le\mathrm C|y|^2|z|.
\]
For the second inequality, subtract the zero perpendicular field on
\(a\mathbf e\) and integrate its derivative along the transverse
segment. The derivative differs from its value at zero by
\(O(|y|^2)\) because the field is odd.
Writing \(w=|z|/|a|\), the upper right derivatives of norms give
\begin{equation}
 \begin{aligned}
 \frac{d|a|}{d\mathfrak s}
 &\ge |a|[\alpha-(\mathrm Cr^2+\kappa_\varepsilon)(1+w)]
                -b_\varepsilon,\\
 D^+w
 &\le-[\gamma-\mathrm Cr^2
             -(\mathrm Cr^2+\kappa_\varepsilon)(1+w)]w
        +(1+w)(\kappa_\varepsilon+b_\varepsilon/|a|).
 \end{aligned}
\label{eq:phil-selection-cone}
\end{equation}
Choose \(W\) larger than twice the initial limiting ratio, then choose
\(r>0\) so small that the terms \(\mathrm Cr^2(1+W)\) consume less than
one quarter of both positive rates. For small enough \(\varepsilon\),
\cref{eq:phil-selection-relative-error} controls the remaining terms.
A bootstrap then keeps \(w\le W\), makes \(|a|\) increase without changing
sign, and gives
\[
 w(\mathfrak s)
 \le w(\mathfrak s_1)e^{-\gamma(\mathfrak s-\mathfrak s_1)/2}
  +\frac{2(1+W)}\gamma
       (\kappa_\varepsilon+b_\varepsilon/|a(t_1)|).
\]
The lower growth bound for \(|a|\) yields an exit time no larger than
\[
 \frac{1+O(r^2)+o(1)}\alpha\log(1/\varepsilon)
\]
in slow time. The corresponding upper norm-growth bound forces the exit
time to tend to infinity. Hence \(w\to0\) at exit, and the exit point
tends to \(\operatorname{sign}(\mathsf j_{\mathrm T})r\mathbf e\).

The reduced diagonal orbit from that point enters any fixed neighborhood
of its matching vertex in finite slow time. The tracking error in
\cref{eq:phil-selection-global-tracking} tends to zero after \(t_1\);
finite-time Gr\"onwall comparison therefore carries the evolving rows into
the same neighborhood. At its selected vertex let
\(\zeta_i=|p_i-p_i^M|\). By the derivative computation above and continuity,
a sufficiently small neighborhood has
\((1-\zeta_i)|\mathfrak g_i(p)|\ge\nu/2\), where \(\nu=1-\mathsf r_E>0\), and
the fitness has the sign directed toward that vertex. Thus
\[
 D^+\zeta_i\le-\tfrac\nu2\zeta_i+
                   \operatorname{TV}(q_t,\mu_{p(t)})
\]
in slow time. The neighborhood remains invariant once the tracking floor
is small enough, and \(\zeta_i\) decays exponentially to a vanishing
floor.

Choose any \(\mathrm C_{\mathrm T}>1/\alpha\), independently of the mean,
and then choose the local radius so that its escape coefficient is less
than \(\mathrm C_{\mathrm T}\). The time remaining after the finite
transit tends to infinity uniformly on the slow-time window
\([\mathrm C_{\mathrm T}\log(1/\varepsilon),
  2\mathrm C_{\mathrm T}\log(1/\varepsilon)]\).
The row distances therefore vanish uniformly on that window. Every
normalized sheet weight is at most the matching weight one, so
\[
 \mu_p(M)\ge\pi_M\ge1-\sum_i\zeta_i,
 \qquad
 q_t(S\ne M)\le\sum_i\zeta_i+
                     \operatorname{TV}(q_t,\mu_{p(t)}).
\]
The last total variation tends to zero by
\cref{eq:phil-selection-global-tracking}; its integrated contraction
exponent at the beginning of the window grows at least as a positive
constant times \(\varepsilon^{-1}\).
This proves the first conclusion for every mean outside the null set in
\cref{eq:phil-selection-nonvanishing}. For the frozen arm,
\cref{eq:phil-selection-harris} and rotation symmetry prove the second
conclusion at every \(\mathrm T>0\) and every mean. The convergence is
pointwise in the fixed mean; no uniform bound near the zero set of the
imprint has been used. \(\square\)

\section{Proof of Joint Selection Along a Schedule}
\label{sec:phil-joint-selection-proof}

\noindent\emph{Proof of \Cref{cor:phil-joint-selection}.}
The saddle estimate \cref{eq:phil-selection-saddle-spectrum} gives
\(\alpha(\mathrm T)\ge7/5\) throughout the stated temperature range.
The proof of \Cref{conj:phil-contrast-formation} permits any window
constant greater than \(1/\alpha(\mathrm T)\), independently of the
initial mean. Thus the common choice \(2\) is valid. Write
\[
 \tau_\varepsilon=2\varepsilon^{-1}\log(1/\varepsilon),
 \qquad W_\varepsilon=[\tau_\varepsilon,2\tau_\varepsilon],
 \qquad 0<\varepsilon<e^{-1}.
\]
For a perfect matching \(M\), define
\[
 \nu_{\mathrm T}^M
 =\mathcal N(0,\mathrm T\mathsf A_M^{-1})\otimes\delta_M,
 \qquad
 \mu_{\mathrm T}=\mu_{\mathrm T,1/2}.
\]
Here \(\mu_{\mathrm T,p}\) is the frozen-row law used in the preceding
proof, with its temperature now displayed explicitly. The inverse
exists because \(\mathsf A_M\succeq\mathrm TI\). Let
\(M_{\mathrm T}(\mathsf m_0)\) be the matching selected by the sign of
\(\mathsf j_{\mathrm T}(\mathsf m_0)\), choosing either matching on its
null zero set.

First record the full one-time consequence of the fixed-temperature
proof. Conditional on assignment \(M\), the law
\(\mu_{\mathrm T,p}\) has the Gaussian factor of \(\nu_{\mathrm T}^M\),
independently of \(p\). The normalized sheet weights are at most the
matching weight one. For
\(\zeta_i=|p_i-p_i^M|\), therefore,
\[
 \mu_{\mathrm T,p}(M)\ge\pi_M\ge1-\sum_i\zeta_i,
 \qquad
 \operatorname{TV}(\mu_{\mathrm T,p},\nu_{\mathrm T}^M)
 =1-\mu_{\mathrm T,p}(M)\le\sum_i\zeta_i.
\]
The row convergence and
\cref{eq:phil-selection-global-tracking} imply
\[
 \sup_{t\in W_\varepsilon}
 \operatorname{TV}
 (q_t^{\mathrm T,\varepsilon,\mathsf m_0},
       \nu_{\mathrm T}^{M_{\mathrm T}(\mathsf m_0)})
 \longrightarrow0
\]
at fixed temperature outside the imprint's null zero set.
Moreover \cref{eq:phil-selection-tracking} itself gives convergence
of the learned relative entropy to zero, uniformly on this window.
For frozen rows, the entropy identity and contraction give
\[
 \KL(q_t^{\mathrm T,0,\mathsf m_0}\Vert\mu_{\mathrm T})
 \le e^{-\lambda(\mathrm T,1/2)t}
       \KL(q_0^{\mathrm T,0,\mathsf m_0}\Vert\mu_{\mathrm T}).
\]
The initial entropy is finite by the Gaussian density bounds, and
\(\lambda(\mathrm T,1/2)>0\). Frozen entropy and total variation
therefore tend to zero on the window for every fixed mean.

Use the explicit entropy majorants as follows. Let
\[
 G_{\mathrm T}=\sup_{p\in[0,1]^4}\|\mathfrak g(p)\|_1,\qquad
 \ell_{\mathrm T}(\varepsilon)
 =\frac{\mathrm c_{\mathrm T}}
 {\mathrm K_{\mathrm T}+8\log2+32\log(1/\varepsilon)},\qquad
 \lambda_u(\mathrm T)
 =\frac{\mathrm c_{\mathrm T}}{\mathrm K_{\mathrm T}+8\log2}.
\]
The first constant is finite. The row lower bound and
\cref{eq:phil-selection-entropy-contraction} make
\(\ell_{\mathrm T}(\varepsilon)\) a rate lower bound throughout
\(t\le2\tau_\varepsilon\), and \(\lambda_u(\mathrm T)\) a rate lower
bound at uniform rows. The common initial entropy equals
\[
 H_0(\mathrm T,\mathsf m_0)
 =\log\mathsf Z(1/2)-6+
 \frac{\operatorname{tr}\overline{\mathsf A}
              +\mathsf m_0^\top\overline{\mathsf A}\mathsf m_0}
      {2\mathrm T},
 \qquad
 \overline{\mathsf A}=\frac1{16}\sum_\sigma\mathsf A_\sigma.
\]
Indeed the log density ratio is
\(\log\mathsf Z(1/2)-|\mathbf x-\mathsf m_0|^2/2+
\mathbf x^\top\mathsf A_\sigma\mathbf x/(2\mathrm T)\);
taking its initial expectation proves the formula.
Thus \(H_0\) is continuous and nonnegative. Uniformly on
\(W_\varepsilon\), the entropy tracking estimate gives
\[
 \begin{aligned}
 \KL(q_t^{\mathrm T,\varepsilon,\mathsf m_0}
               \Vert\mu_{\mathrm T,p(t)})
 &\le B_L(\mathrm T,\varepsilon,\mathsf m_0)\\
 &:=\left[
 e^{-\ell_{\mathrm T}(\varepsilon)\tau_\varepsilon/2}\sqrt{H_0}
 +\frac{\varepsilon G_{\mathrm T}}
             {\sqrt2\ell_{\mathrm T}(\varepsilon)}
 \right]^2,\\
 \KL(q_t^{\mathrm T,0,\mathsf m_0}\Vert\mu_{\mathrm T})
 &\le B_F(\mathrm T,\varepsilon,\mathsf m_0)
 :=e^{-\lambda_u(\mathrm T)\tau_\varepsilon}H_0.
 \end{aligned}
\]
These majorants are continuous in positive mobility and mean and tend
to zero at fixed temperature and mean as \(\varepsilon\downarrow0\).

Let \(\mathcal E_{\mathrm T}(\varepsilon,\mathsf m_0)\) be the minimum
of one and the maximum of the following five nonnegative quantities:
\[
 \begin{aligned}
 &\sup_{t\in W_\varepsilon}
          \max_i|p_i(t)-p_i^{M_{\mathrm T}(\mathsf m_0)}|,\\
 &\sup_{t\in W_\varepsilon}
       \operatorname{TV}
       (q_t^{\mathrm T,\varepsilon,\mathsf m_0},
                   \nu_{\mathrm T}^{M_{\mathrm T}(\mathsf m_0)}),\\
 &B_L(\mathrm T,\varepsilon,\mathsf m_0),\\
 &\sup_{t\in W_\varepsilon}
       \operatorname{TV}
          (q_t^{\mathrm T,0,\mathsf m_0},\mu_{\mathrm T}),\\
 &B_F(\mathrm T,\varepsilon,\mathsf m_0).
 \end{aligned}
\]
The preceding estimates prove
\(\mathcal E_{\mathrm T}(\varepsilon,\mathsf m_0)\to0\)
for almost every fixed mean at each fixed temperature.

We will take a supremum over mobility, so measurability must be checked.
Fix a temperature and restrict the positive mobility and the mean to
compact sets, with the mobility bounded away from zero. The corresponding
horizons are bounded, and
\cref{eq:phil-selection-row-bounds} bounds every row away from zero and
one. On this row set the posterior redraw probabilities are uniformly
Lipschitz in the rows, independently of the configuration.
Maximally couple the initial laws and then use common Brownian motions
and proposal clocks until the first different redraw. The law difference
is bounded by the initial total variation plus a constant times the
integral of row differences. Subtracting the learning equations bounds
the row difference by integrals of row and law differences and by the
mobility difference. Gr\"onwall's inequality proves total-variation
continuity in the initial mean and positive mobility on each bounded
time interval. Time continuity follows from
\cref{eq:phil-selection-forward}.

Parametrizing the window by \(t=u\tau_\varepsilon\),
\(1\le u\le2\), proves continuity in positive mobility of each row or
total-variation window supremum. The explicit entropy majorants are
continuous by their formulas. The matching as a function of the mean is
Borel because its imprint is analytic. Thus the error is measurable
in the mean, and its supremum over an interval of positive mobilities
equals its supremum over rational mobilities. No uniform continuity at
zero mobility or separate entropy-continuity assertion is needed.

Fix the sequence \(\mathrm T_n\downarrow0\) in the stated temperature
interval and \(0<\eta_n<1\) with \(\eta_n\to0\). For the prescribed
exponent \(a>7/2\), put \(h_n=\mathrm T_n^a\). The following argument
in fact permits arbitrary positive caps \(h_n\).
On the mean space \(\R^{12}\), introduce the auxiliary probability
measure \(\rho=\mathcal N(0,I_{12})\). Its density is strictly positive;
it introduces no assumption that the physical mean is random.
Set
\[
 F_n(r,\mathsf m_0)
 =\sup_{0<\varepsilon\le r}
       \mathcal E_{\mathrm T_n}(\varepsilon,\mathsf m_0).
\]
The measurable functions \(F_n\) are bounded by one. At fixed \(n\),
they tend to zero for \(\rho\)-almost every mean as \(r\downarrow0\).
Dominated convergence yields
\[
 \int F_n(r,\mathsf m_0)d\rho(\mathsf m_0)\longrightarrow0.
\]
Choose \(d_n=2^{-k_n}\) recursively, where \(k_n\) is the smallest
positive integer satisfying
\[
 d_n\le\min(e^{-1},2^{-n},h_n,d_{n-1}/2),
 \qquad
 \int F_n(d_n,\mathsf m_0)d\rho(\mathsf m_0)
       \le\eta_n2^{-n},
\]
with the \(d_{n-1}/2\) condition omitted for \(n=1\). The preceding
limit proves existence. This defines deterministic thresholds
independent of the mean; it does not give numerical values or an explicit
power law.

Let \(B_n=\{\mathsf m_0:F_n(d_n,\mathsf m_0)>\eta_n\}\).
Markov's inequality and countable subadditivity give
\[
 \rho(B_n)\le2^{-n},
 \qquad
 \rho\left(\bigcup_{n\ge N}B_n\right)\le2^{1-N}.
\]
The decreasing intersection of these tail unions consequently has
measure zero. Add the countable union of the imprint zero sets. The
resulting exceptional set is Lebesgue-null because \(\rho\) and Lebesgue
measure have the same null sets. Outside it, for every sufficiently large
\(n\), all \(0<\varepsilon\le d_n\) satisfy
\(F_n(d_n,\mathsf m_0)\le\eta_n\). Since \(\eta_n<1\), clipping has
no effect on this inequality: the row and total-variation errors and both entropy majorants are at
most \(\eta_n\). The majorants bound the actual entropies, giving
all five stated conclusions.

The frozen matching probability is, exactly,
\[
 \mu_{\mathrm T}(M)
 =\frac1{2+8\mathsf r_E+4\mathsf r_N+2\mathsf r_C}
 =\frac1{2D_0(\mathrm T)}
 =\frac12+O(\mathrm T).
\]
Indeed \(\mathsf r_E,\mathsf r_N=O(\mathrm T)\) and
\(\mathsf r_C=O(\mathrm T^{3/2})\), as proved from the sheet
determinants. Pushforward cannot increase total variation. If
\(q^{S}\) denotes the assignment marginal, the two established
total-variation estimates therefore imply, uniformly on \(W_\varepsilon\),
\[
 \left|
 \operatorname{TV}
 (q_t^{\mathrm T_n,\varepsilon,\mathsf m_0,S},
                  q_t^{\mathrm T_n,0,\mathsf m_0,S})
 -\left(1-\frac1{2D_0(\mathrm T_n)}\right)
 \right|\le2\eta_n.
\]
Here the comparison distance is exactly
\(\operatorname{TV}(\delta_M,\mu_{\mathrm T_n}^S)
=1-\mu_{\mathrm T_n}(M)\), for either matching.
The assignment-law contrast thus tends to \(1/2\), even if the selected
matching changes with \(n\).

Finally, defining \(\mathrm T(\varepsilon)=\mathrm T_n\) on
\(d_{n+1}<\varepsilon\le d_n\) gives a single stepwise path with both
parameters tending to zero. The proved bounds hold for all sufficiently
small mobilities on this path. Arbitrary mobility caps are admissible;
for example \(h_n=\mathrm T_n^n\) forces
\(\varepsilon_n/\mathrm T_n^a\to0\) for each fixed \(a>0\) whenever
\(\varepsilon_n\le d_n\). This imposes an upper bound and does not certify
a prescribed equality between temperature and mobility.

Nothing in this argument makes the imprint sign eventually constant as
temperature vanishes. Its fixed-temperature nonvanishing and null zero
set give no quantitative lower bound at a fixed mean in that limit.
Accordingly the matching is allowed to depend on \(n\), and the result
concerns complete one-time laws, not complete histories or conditional
first-formation events. \(\square\)





\subsection*{An explicit temperature bound for tracking}

\propositionheading{Temperature-dependent entropy tracking}{prop:phil-temperature-tracking}
In the four-cycle of \cref{conj:phil-contrast-formation}, let
\(0<\mathrm T\le1/100\), \(p\in(0,1)^4\), and
\(p_{\min}=\min_i\{p_i,1-p_i\}\). Define
\[
 \mathrm K_0=210\sqrt{2\pi/5},\qquad
 \mathrm c_* =\min\left\{
 \frac2{1+4\mathrm K_0e^{13/8}},\frac1{\mathrm K_0e^{13/4}}\right\}.
\]
The entropy production of the full frozen law satisfies
\[
 \mathscr D_p(f)\ge
 \frac{\mathrm c_*\mathrm T^{5/2}}
 {4+3\log(1/\mathrm T)+8\log(1/p_{\min})}
 \operatorname{Ent}_{\mu_p}(f).
\]
Write \(\varepsilon=\epsilon_{\mathrm L}\),
\(t_{\mathrm L}=2\varepsilon^{-1}\log(1/\varepsilon)\), and
\[
 \mathrm D(\mathrm T,\varepsilon)
 =4+3\log(1/\mathrm T)+8\log2+32\log(1/\varepsilon),\qquad
 \ell=\frac{\mathrm c_*\mathrm T^{5/2}}
                {\mathrm D(\mathrm T,\varepsilon)}.
\]
For \(0<\varepsilon<e^{-1}\) and \(0\le t\le2t_{\mathrm L}\),
\begin{equation}
 \operatorname{TV}(q_t,\mu_{p(t)})
 \le e^{-\ell t/2}
       \sqrt{\frac{13(12+|\mathsf m_0|^2)}{4\mathrm T}}
    +\frac{2\varepsilon\mathrm D(\mathrm T,\varepsilon)}
                 {\mathrm c_*\mathrm T^{7/2}}.
\label{eq:phil-temperature-tracking}
\end{equation}
In particular, for \(\varepsilon=\mathrm T^a\) with any \(a>7/2\),
this total variation tends to zero uniformly on
\([t_{\mathrm L},2t_{\mathrm L}]\) and on bounded sets of initial means.
This conclusion compares one-time laws; it supplies no imprint lower bound.
\status{ESTABLISHED}

\paragraph{Proof.}
Use \(u=(\bar k+\bar m)/\sqrt2\),
\(v=(\bar k-\bar m)/\sqrt2\), and \(\omega\), and put
\(R=[D,-\sqrt2I]\). The common kernel
\(K=\ker R\oplus\ker D\) has dimension five. Every
\(\mathsf H_\sigma\) vanishes there, and the reference precision preserves
\(K\) and \(N=K^\perp\). Its restrictions have eigenvalues one and five.
The passive process on \(K\) is therefore an independent Gaussian
Ornstein--Uhlenbeck process with entropy contraction at least \(2\mathrm T\).
The active process on \(N\), together with the assignment, is autonomous:
all redraw costs depend only on its residuals. Rescaling its coordinate by
\(\mathrm T^{-1/2}\) gives unit diffusion and drift matrix
\(a_\sigma=(\mathsf H_\sigma+\mathrm T\mathsf\Omega)|_N\).
Since every edge multiplicity is at most two,
\[
 \mathrm TI\preceq a_\sigma\preceq13I.
\]
Indeed \(RR^\top=DD^\top+2I\) has eigenvalues \(2,4,4,6\),
while \(D^\top D\) has eigenvalues \(0,2,2,4\), so
\(\|\mathsf H_\sigma\|\le12\).

Divide the four determinants in
\cref{eq:phil-selection-sheet-determinants} by the passive determinant
\(5\mathrm T^5\). Their positive coefficients, evaluated at
\(\mathrm T=1/100\) for upper bounds, give the active bounds
\[
 \begin{aligned}
 5120\mathrm T^3&\le d_M\le5262\mathrm T^3,&
 896\mathrm T&\le d_E\le1017\mathrm T,\\
 896\mathrm T&\le d_N\le995\mathrm T,&
 3072&\le d_C<5120.
 \end{aligned}
\]
Thus \(w_M\) and \(w_C\) are the largest and smallest active sheet
weights, and \(w_M\le\mathrm T^{-3/2}/\sqrt{5120}\),
\(w_C/w_M\ge\mathrm T^{3/2}\). The finite-label constant in the
entropy proof consequently satisfies
\[
 L_p\le4+3\log(1/\mathrm T)+8\log(1/p_{\min}).
\]

Apply the Gaussian-sheet and prior-weighted-ball argument of
\cref{eq:phil-selection-gaussian-inequalities,eq:phil-selection-entropy-contraction}
in the seven active coordinates, using the ball of radius \(1/2\), whose
volume is \(\pi^3/840\). Its Gaussian comparison factor is at most
\(4\mathrm K_0e^{13/8}\mathrm T^{-3/2}\), because
\[
 \frac{(2\pi)^{7/2}}{(\pi^3/840)\sqrt{5120}}=\mathrm K_0.
\]
More explicitly, the two energy coefficients in that argument obey
\[
 C_c\le(2+8\mathrm K_0e^{13/8})L_p\mathrm T^{-5/2},\qquad
 C_j\le4\mathrm K_0e^{13/4}L_p\mathrm T^{-5/2}.
\]
The continuous production equals four times its quadratic form and the
jump production bounds four times its quadratic form. Hence the active
rate is at least \(\mathrm c_*\mathrm T^{5/2}/L_p\). Entropy decomposes
over the independent passive and active factors. Cauchy--Schwarz bounds
the energy of a marginal by the corresponding integrated energy, so the
full rate is at least the smaller of the two rates. The displayed active
lower bound is less than \(2\mathrm T\), proving the full-state inequality.
The form-domain extension uses the positive regularization and cutoffs
already used for \cref{eq:phil-selection-entropy-contraction}.

A one-receiver flip joins only the class pairs \(M\)--\(E\),
\(N\)--\(E\), and \(C\)--\(E\). The determinant bounds above make every
ratio of such adjacent weights at most \(1/\mathrm T\).
Writing the evidence as \(p_iZ_i^++(1-p_i)Z_i^-\), the ratio
\(Z_i^+/Z_i^-\) is between the extremal adjacent ratios. Therefore
\(|\partial_{p_i}\log Z|\le1/\mathrm T\), and
\(\|\mathfrak g\|_1\le4/\mathrm T\).
The row lower bound gives the rate floor \(\ell\) on the stated horizon.
The initial entropy is at most
\(13(12+|\mathsf m_0|^2)/(2\mathrm T)\): in the Gaussian density ratio,
\(\mathsf B_\sigma\preceq13I/\mathrm T\),
\(\mathsf Z\le1\), and the remaining quadratic term has nonpositive
expectation. Insert these bounds in
\cref{eq:phil-selection-tracking} and use Pinsker's inequality to obtain
\cref{eq:phil-temperature-tracking}.
For \(\varepsilon=\mathrm T^a\), its second term is
\(O(\mathrm T^{a-7/2}\log(1/\mathrm T))\), while
\(\ell t_{\mathrm L}\) grows as a positive constant times
\(\mathrm T^{5/2-a}\). The exponential therefore dominates the initial
\(\mathrm T^{-1/2}\) factor, proving the last assertion. \(\square\)
\section{Coherence and Persistence of the Selected Pairing}
\label{sec:phil-selected-coherence-proof}

\propositionheading{Coherence and persistence after pairing selection}{prop:phil-selected-pair-persistence}
In the four-cycle of \cref{conj:phil-contrast-formation}, let
\(0<\mathrm T\le1/100\), let \(M\) be one matching, and suppose
\([t_0,t_0+\mathsf h]\subseteq[t_{\mathrm L},2t_{\mathrm L}]\), with
\(\mathsf h>0\). Write
\(\nu_{\mathrm T}^M=\mathcal N(0,\mathrm T\mathsf A_M^{-1})\otimes\delta_M\).
Assume the learned and frozen recognition laws satisfy
\[
 \sup_{t_0\le t\le t_0+\mathsf h}
 \operatorname{TV}(q_t^{\mathrm L},\nu_{\mathrm T}^M)\le\delta,
 \qquad
 \operatorname{TV}(q_{t_0}^{0},\mu_{\mathrm T,1/2})\le\delta.
\]
Let \(\mathsf O_{\mathrm T,\mathsf h}^M\) be the stationary
Ornstein--Uhlenbeck path law with drift \(-\mathsf A_M\mathbf x\),
diffusion covariance \(2\mathrm TI\), and constant assignment \(M\).
For the configuration and assignment path laws on this interval,
\begin{align}
 \operatorname{TV}(\mathsf P^{\mathrm L}_{t_0,\mathsf h},
                    \mathsf O_{\mathrm T,\mathsf h}^M)
 &\le(\mathsf h+2)\delta,
 \label{eq:phil-selected-path-coupling}\\
 \operatorname{TV}\left(\mathsf P^{0}_{t_0,\mathsf h},
   \tfrac12\mathsf O_{\mathrm T,\mathsf h}^{M_0}
        +\tfrac12\mathsf O_{\mathrm T,\mathsf h}^{M_1}\right)
 &\le\delta+16\mathrm T+31\mathrm T\mathsf h.
 \label{eq:phil-frozen-path-mixture}
\end{align}
The frozen arm keeps its original uniform rows.

For resolution \(\mathsf a>0\), let
\(\mathscr E_M(\mathsf a,\mathsf h)\) require throughout the interval
that the assignment is \(M\), the stacked residual norm
\(\sqrt{(r^b)^2+(r^m)^2}\) on its two edges is at most
\(\mathsf a\), and that norm on each exterior edge exceeds
\(\mathsf a\).
Define
\[
 \mathsf b(\mathrm T,\mathsf a,\mathsf h)
 =8\sqrt{2/\pi}\mathsf a+
            (8+3400\mathsf h)\mathrm T/\mathsf a^2.
\]
Then
\begin{equation}
 \Pr_{\mathrm L}(\mathscr E_M^c)
 \le(\mathsf h+2)\delta+\mathsf b(\mathrm T,\mathsf a,\mathsf h).
 \label{eq:phil-selected-coherence-persistence}
\end{equation}
In particular, the choice
\(\mathsf a=\mathrm T^{1/4}\),
\(\mathsf h=\mathrm T^{-1/4}\), and \(\delta\le\mathrm T\)
makes this probability tend to zero.
The matching-window label, equal to \(M_j\) when
\(\mathscr E_{M_j}\) holds and to an additional symbol otherwise,
converges to \(\delta_M\) in the learned arm and to the equal matching
mixture in the frozen arm.

At one time, define \(\mathfrak P_{\mathsf a}(\mathbf x)\) as the
connected-component partition of edges whose stacked residual norm
\(\sqrt{(r^b)^2+(r^m)^2}\) is at most \(\mathsf a\).
For times when both
one-time hypotheses above hold, put
\[
 \mathsf e(\mathrm T,\mathsf a)
 =4\sqrt{2/\pi}\mathsf a+4\mathrm T/\mathsf a^2.
\]
The learned and frozen partition laws satisfy
\begin{equation}
 \left|\operatorname{TV}
   (q_t^{\mathrm L}\mathfrak P_{\mathsf a},
                 q_t^0\mathfrak P_{\mathsf a})-\tfrac12\right|
 \le2\delta+16\mathrm T+2\mathsf e(\mathrm T,\mathsf a).
 \label{eq:phil-coherent-partition-contrast}
\end{equation}
Thus the contrast persists for the coherent partition, not only the
source label.

These conclusions apply along \cref{cor:phil-joint-selection} with
tolerance \(\delta\le\mathrm T\), mobility cap
\(\epsilon_{\mathrm L}\le\mathrm T^2\), and \(t_0=t_{\mathrm L}\).
The matching can vary along that schedule. The internal restoring rates
are \(8+\mathrm T\) and \(4+5\mathrm T\), on an event whose probability
tends to one. The resolution therefore has restoring ratio at least
\(4\mathsf a^2/\mathrm T\to\infty\).
If the rows are also within \(\delta\) of their matching vertex,
the stationary law with those rows frozen realizes the same persistence
event with probability at least
\(1-4(\mathsf h+2)\delta-\mathsf b\).
This late coherence is restored equilibrium coherence; it supplies no
entrance or first-formation assertion.
\status{ESTABLISHED}

\noindent\emph{Proof.}
First retain the proposal mechanism in the comparison. Put
\(r(t)=q_t^{\mathrm L}(S\ne M)\), and let
\(b_{\ne M}(t,\mathbf x)\) be the probability of proposing a different
assignment at a clock ring. Since the clock has rate one,
\[
 r'(t)=\E b_{\ne M}(t,\mathbf X_t)-r(t).
\]
If \(N_{\ne M}\) counts proposals to an assignment other than \(M\)
on the interval, its compensator gives
\begin{equation}
 \E N_{\ne M}
 =r(t_0+\mathsf h)-r(t_0)+
                   \int_{t_0}^{t_0+\mathsf h}r(t)dt.
 \label{eq:phil-selected-proposal-occupation}
\end{equation}
Consequently
\(\Pr(\exists t\in[t_0,t_0+\mathsf h]:S_t\ne M)
 \le(\mathsf h+1)\delta\).
Couple the initial law maximally with \(\nu_{\mathrm T}^M\), and use
the same future Brownian motion for the native process and the
stationary matching process. Their paths agree until an initial
mismatch or a proposal to another assignment. The union bound and
\cref{eq:phil-selected-proposal-occupation} prove
\cref{eq:phil-selected-path-coupling}. This argument counts proposals
of the native process and does not condition its diffusion on an
assignment history.

For the frozen comparison, use the normalized weights
\(1,\mathsf r_E,\mathsf r_N,\mathsf r_C\) from
\cref{sec:phil-pairing-selection-proof}. The determinant bounds in the
proof of \cref{prop:phil-temperature-tracking} give
\[
 \mathsf r_E,\mathsf r_N\le\tfrac52\mathrm T,
 \qquad \mathsf r_C\le\tfrac32\mathrm T^{3/2}.
\]
The invariant mass \(\kappa_{\mathrm T}\) of nonmatching assignments
is therefore at most \(16\mathrm T\). Conditional on either matching,
its configuration law is the Gaussian factor of \(\nu_{\mathrm T}^M\).
The invariant law is consequently within \(\kappa_{\mathrm T}\) of
the equal mixture of these two laws.

Detailed balance bounds the stationary flux from \(M\) to all
nonmatching assignments by \(\kappa_{\mathrm T}\).
For the other matching \(M'\), write \(\mu_M(\mathbf x)\) for the
joint invariant density on the matching sheet. Pointwise,
\[
 \mu_M b_{M'}
 =\frac{\mu_M\mu_{M'}}{\sum_S\mu_S}
 \le\tfrac12\sqrt{\mu_M\mu_{M'}}.
\]
Since \((\mathsf A_M+\mathsf A_{M'})/2\) is the oriented-cycle
matrix, Gaussian integration gives
\(\int\sqrt{\mu_M\mu_{M'}}=\mu_{\mathrm T,1/2}(M)\mathsf r_C\).
Dividing the two flux bounds by the matching mass yields
\[
 \E_{\nu_{\mathrm T}^M}b_{\ne M}
 \le8\mathsf r_E+4\mathsf r_N+\tfrac52\mathsf r_C
 \le31\mathrm T.
\]
Thus the stationary flux out of the two matching sheets together is
at most \(31\mathrm T\). Couple the stationary native process at
its initial time with the equal matching mixture and use common
Brownian increments until the first departure from its coupled
matching. Its failure probability is at most
\(16\mathrm T+31\mathrm T\mathsf h\). The frozen native process
started at \(q_{t_0}^0\) differs from its stationary path law by at
most its initial total variation \(\delta\), because the frozen
transition kernel is the same. This proves
\cref{eq:phil-frozen-path-mixture}.

It remains to bound residuals on the comparison path.
The four normal coordinates
\(z_g^\pm=n_g^\pm\cdot\mathbf X\), \(g\in M\), from
\cref{eq:rg-formation-cycle-normals} have rates
\(8+\mathrm T\) and \(4+5\mathrm T\). Their stationary variances
are at most \(\mathrm T/4\), and
\[
 r_g^b=\sqrt2z_g^++z_g^-,\qquad
 r_g^m=\sqrt2z_g^+-z_g^-.
\]
For a scalar Ornstein--Uhlenbeck noise starting at zero, integration
by parts in its stochastic convolution gives
\[
 Z_t=\sqrt{2\mathrm T}
 \left[W_t-\lambda\int_0^te^{-\lambda(t-u)}W_u du\right],
 \qquad
 \E\sup_{t\le\mathsf h}|Z_t|^2\le32\mathrm T\mathsf h.
\]
The last inequality uses
\(\sup|Z_t|\le2\sqrt{2\mathrm T}\sup|W_t|\) and the Brownian
maximal estimate proved in the matching-episode argument of
\cref{sec:rg-formation-cycle-proof}.
For the four-dimensional stationary normal vector \(z\), it follows
that
\[
 \E\sup_{t\le\mathsf h}|z_t|^2
 \le2\mathrm T+256\mathrm T\mathsf h.
\]
Each selected stacked residual norm is at most \(2|z_t|\), since
\((r_g^b)^2+(r_g^m)^2=4(z_g^+)^2+2(z_g^-)^2\).
Thus the probability that any internal stacked norm exceeds
\(\mathsf a\) at any time is at most
\((8+1024\mathsf h)\mathrm T/\mathsf a^2\).

Every exterior link coordinate is absent from \(H_M\), and its
block in \(\mathsf A_M\) is the scalar \(\mathrm T\).
It is therefore an independent \(\mathcal N(0,1)\) coordinate in
the stationary Gaussian. Each exterior belief or model residual
has variance at least one. The probability that any of these four
residuals initially has absolute value at most \(2\mathsf a\)
is at most \(8\sqrt{2/\pi}\mathsf a\).

To control their increments, diagonalize the symmetric positive
matrix \(\mathsf A_M\). In a scalar eigencoordinate with rate
\(\lambda>0\), stationarity and the preceding noise estimate imply
\[
 \E\sup_{t\le\mathsf h}|X_t-X_0|^2
 \le\frac{2\mathrm T}{\lambda}(1-e^{-\lambda\mathsf h})^2
                       +64\mathrm T\mathsf h
 \le66\mathrm T\mathsf h,
\]
since \((1-e^{-v})^2\le v\) for \(v\ge0\).
Summing twelve coordinates gives
\begin{equation}
 \E\sup_{t\le\mathsf h}|\mathbf X_t-\mathbf X_0|^2
 \le792\mathrm T\mathsf h.
 \label{eq:phil-selected-ou-increments}
\end{equation}
Every residual covector has norm \(\sqrt3\). Except with
probability \(2376\mathrm T\mathsf h/\mathsf a^2\), all exterior
increments are therefore at most \(\mathsf a\) in absolute value.
The initial small-ball estimate and the two maximal estimates prove
\(\mathsf O_{\mathrm T,\mathsf h}^M(\mathscr E_M^c)\le\mathsf b\).
Together with \cref{eq:phil-selected-path-coupling}, this proves
\cref{eq:phil-selected-coherence-persistence}.
The frozen path comparison proves the stated matching-window label
limit whenever \(\mathsf b\to0\),
\(\mathsf h\delta\to0\), and \(\mathrm T\mathsf h\to0\), in
particular for the displayed choices.

For a single time no increment estimate is needed. The normal second
moment is at most \(\mathrm T\), so the internal failure probability
is at most \(4\mathrm T/\mathsf a^2\). The exterior small-ball
bound at radius \(\mathsf a\) is
\(4\sqrt{2/\pi}\mathsf a\). Thus the partition induced by
\(\nu_{\mathrm T}^{M_j}\) differs from the matching partition
\(\mathfrak P_j\) with probability at most \(\mathsf e\).
The learned partition law is within \(\delta+\mathsf e\) of
\(\delta_{\mathfrak P_M}\); the frozen one is within
\(\delta+16\mathrm T+\mathsf e\) of
\((\delta_{\mathfrak P_0}+\delta_{\mathfrak P_1})/2\).
These matching partitions are distinct, so the triangle inequality
for total variation proves \cref{eq:phil-coherent-partition-contrast}.

The schedule in \cref{cor:phil-joint-selection} permits tolerance
\(\mathrm T\) and mobility cap \(\mathrm T^2\).
Then \(\mathrm T^{-1/4}\le t_{\mathrm L}\) for sufficiently small
\(\mathrm T\), so the following persistence window lies within the
selection window. All errors just bounded tend to zero.
The normal rates establish the restoring ratio on the native event.
Finally, if \(\max_i|p_i-p_i^M|\le\delta\),
\(\operatorname{TV}(\mu_{\mathrm T,p},\nu_{\mathrm T}^M)
 =1-\mu_{\mathrm T,p}(M)\le4\delta\).
Applying the same occupation and coupling argument to this
stationary frozen-row process gives the stationary lower bound in
the statement. The persistence event is therefore already typical
at equilibrium for those rows and carries no first-entrance claim.
\hfill\(\square\)



\subsection{The first contrast on the original episode window}

\propositionheading{A first-block null under slower source learning}{prop:phil-early-block-null}
In the model of \cref{conj:phil-contrast-formation}, let
\(0<\mathrm T\le10^{-16}\), and let the initial law be
\(\mathcal N(\mathsf m_0,I_{12})\) times uniform assignments.
The learned and frozen arms have the same initial mean and differ only in
the source-learning mobility \(\epsilon_{\mathrm L}\).
Set
\[
 H_{\mathrm T}=2\mathrm T^{-1/8}.
\]
Keep the contrast event of \cref{prop:phil-seeded-contrast-formation}
with its regime parameter equal to \(\mathrm T^{1/8}\): an eligible
proper connected block has an internal belief relative entropy exceeding
\(1/4\) initially, and for a duration-one interval inside
\([0,H_{\mathrm T}]\) all its internal belief and model relative
entropies are at most \(\sqrt{\mathrm T}\), while its boundary belief
relative entropies exceed \(1/2\).
For each block take the infimum of its qualifying completion times,
with value infinity when there is no such interval. Let \(K_{\mathrm T}\)
be the block with least such infimum, using one fixed tie rule in both
arms, and give it a cemetery value when no block qualifies. Write
\(E_{\mathrm T}\) for the event that its value is not the cemetery value.
The infimum need not itself be a qualifying completion time.

Define the positive, temperature-independent constant
\[
 \mathsf c(\mathsf m_0)
 =\frac{e^{-1/4}\varrho_\circ^{12}}{737280}
   \exp\left[-\frac12
       (|\mathbf x^\circ-\mathsf m_0|+\varrho_\circ)^2\right],
 \qquad
 \mathsf d_{\mathrm T}
 =\min\{1,8\epsilon_{\mathrm L}\mathrm T^{-1/4}\}.
\]
Then
\[
 \Pr_0(E_{\mathrm T})\ge\mathsf c(\mathsf m_0),
 \qquad
 \Pr_{\epsilon_{\mathrm L}}(E_{\mathrm T})
       \ge\mathsf c(\mathsf m_0)-\mathsf d_{\mathrm T}.
\]
If \(\mathsf d_{\mathrm T}<\mathsf c(\mathsf m_0)\), the two
conditional block laws satisfy
\[
 \operatorname{TV}\left(
 \operatorname{Law}_{\epsilon_{\mathrm L}}(K_{\mathrm T}\mid E_{\mathrm T}),
 \operatorname{Law}_0(K_{\mathrm T}\mid E_{\mathrm T})\right)
 \le\frac{\mathsf d_{\mathrm T}}{\mathsf c(\mathsf m_0)}.
\]
In particular their difference vanishes for every fixed Gaussian mean
whenever
\(\mathrm T\downarrow0\) and
\(\epsilon_{\mathrm L}\mathrm T^{-1/4}\to0\).
This includes the caps
\(\epsilon_{\mathrm L}\le\mathrm T_n^a\), \(a>7/2\), in
\cref{cor:phil-joint-selection}.
\status{ESTABLISHED}

\noindent\emph{Proof.}
First prove a uniform comparison of the complete native path laws on
any finite horizon \(H\). The learned rows are deterministic and satisfy
\[
 |p_i(t)-1/2|\le\epsilon_{\mathrm L}t/2.
\]
At a common configuration, write \(A,B>0\) for a receiver's likelihood
factors for its two sources. Its posterior positive-source probabilities
in the two arms are
\[
 b(p)=\frac{pA}{pA+(1-p)B},
 \qquad b(1/2)=\frac A{A+B}.
\]
Direct subtraction gives
\[
 |b(p)-b(1/2)|
 =\frac{AB|2p-1|}{(pA+(1-p)B)(A+B)}
 \le|2p-1|.
\]
Indeed the denominator is
\(pA^2+AB+(1-p)B^2\ge AB\). This bound is independent of
temperature and the likelihood ratio, without an interior-row restriction.

Use identical initial samples, a common Brownian motion, and the same
rate-one proposal clock in the two arms. Until the first different redraw,
their configurations and assignments coincide: a held-sheet drift and its
noise do not depend on the prior rows. At each ring, couple each receiver's
two Bernoulli laws maximally. The union bound and the row estimate bound
the probability of any different receiver choice at time \(t\) by
\[
 \sum_i|b_i(p_i(t))-b_i(1/2)|
 \le4\epsilon_{\mathrm L}t.
\]
The compensator of the rate-one clock, stopped at the first mismatch,
bounds its probability before \(H\) by
\(\int_0^H4\epsilon_{\mathrm L}t dt=2\epsilon_{\mathrm L}H^2\).
Consequently,
\begin{equation}
 \operatorname{TV}\left(
 \operatorname{Law}_{\epsilon_{\mathrm L}}(\mathbf X_{[0,H]},S_{[0,H]}),
 \operatorname{Law}_0(\mathbf X_{[0,H]},S_{[0,H]})\right)
 \le\min\{1,2\epsilon_{\mathrm L}H^2\}.
\label{eq:phil-early-path-coupling}
\end{equation}
The bound follows because the complete paths are identical unless the
coupling has had a mismatch. Taking \(H=H_{\mathrm T}\) gives
\(\mathsf d_{\mathrm T}\).

The first-block mark is a measurable function of these paths. To see this
directly, for a fixed block let \(A_I(t)\) be the maximum of its internal
belief and model relative entropies on \([t-1,t]\), and let \(B_I(t)\)
be the minimum boundary belief relative entropy there. Configuration
paths are continuous, so these two window functions are continuous.
Replace the strict condition \(B_I(t)>1/2\) by the countable union of
conditions \(B_I(t)\ge1/2+1/k\), \(k\ge1\). For each \(k\) the
qualifying end times form a closed subset of the compact interval
\([1,H_{\mathrm T}]\). Its first time is measurable, as may be checked
by minimizing
\(\max(A_I(t)-\sqrt{\mathrm T},1/2+1/k-B_I(t))\) on each compact
initial subinterval. The infimum over \(k\), the initial eligibility
condition, and the finite comparison over blocks preserve measurability.
This construction implements the strict boundary convention without
asserting that its infimum is attained.

A measurable map of paths does not increase total variation. Let
\(a_I,b_I\) be the probabilities of the non-cemetery marks in the learned
and frozen arms, \(p=\sum_Ia_I\), \(q=\sum_Ib_I\), and \(d\) the
total-variation distance of the two cemetery-augmented mark laws. Then
\[
 2d=\sum_I|a_I-b_I|+|p-q|,
 \qquad d\le\mathsf d_{\mathrm T}.
\]
If \(p\ge q>0\), the triangle inequality gives
\[
 \frac12\sum_I\left|\frac{a_I}{p}-\frac{b_I}{q}\right|
 \le\frac{\sum_I|a_I-b_I|+(p-q)}{2p}
 =\frac d p.
\]
Interchanging the arms proves the bound \(d/\max(p,q)\).
The frozen matching-seed sub-event in
\cref{prop:phil-seeded-contrast-formation} belongs to \(E_{\mathrm T}\)
and has probability at least \(\mathsf c(\mathsf m_0)\).
Thus \(q\ge\mathsf c(\mathsf m_0)\) and
\(p\ge q-d\ge\mathsf c(\mathsf m_0)-\mathsf d_{\mathrm T}\).
The claimed conditional bound follows, and its limiting conclusion is
immediate. \(\square\)

The proposition keeps the original contrast event and uses its seeded
sub-event only as a probability lower bound. That sub-event already has
the restoring interaction and stationary comparison of
\cref{prop:rg-formation-endogenous-cycle}; the proof does not certify
every history in the broader contrast event under the coherence-episode
rule. Any common measurable admissibility restriction retaining the
seeded sub-event has the same comparison.

The result does not cover the original coupled choice
\(\epsilon_{\mathrm L}=\mathrm T^{1/8}\), for which its bound is
vacuous. It also does not compare different initial means or establish a
conditional parent error. On the slower schedules, late assignment
selection can coexist with a vanishing difference between the conditional
laws of the first contrast on \([0,H_{\mathrm T}]\). These are
different observables on different horizons.



\section{Activation from the Actual Conditioned Recognition Law}
\label{sec:phil-actual-conditioned-parent}

\propositionheading{Conditional information after a completed coherence prefix}{prop:phil-actual-conditioned-parent}
In the learned four-cycle, fix \(0<\mathrm T\le1/100\), a selected
matching \(M\), and \(s=t_{\mathrm L}\ge1\). Suppose, throughout
\([s,2s]\),
\[
 \KL(q_t\Vert\mu_{\mathrm T,p(t)})\le\eta,
 \qquad
 \operatorname{TV}(q_t,\nu_{\mathrm T}^M)\le\delta.
\]
Set \(\mathsf a=\mathrm T^{1/4}\), \(\theta=s+1\), and let
\(E=\mathscr E_M(\mathsf a,1)\) be the actual matching and coherence
event of \cref{prop:phil-selected-pair-persistence} on \([s,s+1]\).
Assume
\[
 \rho=3\delta+\mathsf b(\mathrm T,\mathsf a,1)<\tfrac12,
 \qquad
 h_b(r)=-r\log r-(1-r)\log(1-r),
 \qquad
 \kappa_E=\frac{\eta+h_b(\rho)}{1-\rho}.
\]
For all \(t\ge\theta\), write
\(q_t^E=\operatorname{Law}(\mathbf X_t,S_t\mid E)\), always
conditioning on this same completed prefix. Keep the single marked
channel \(C\) of \cref{prop:rg-formation-endogenous-cycle}(b), including
its identity action off the matching strata and its identity extension
on evaluator samples. The parent activation law is exactly
\[
 \mu_{\mathrm{act}}=C_\#q_\theta^E.
\]

Let \(\mathcal H_{\mathrm{opt}}(q)\) be the mean discarded conditional
relative entropy of \(q\) against the restricted optimizer under \(C\).
Let \(\mathcal H_{\mathrm{phys}}(q)\) be the corresponding mean
discarded conditional relative entropy against the full physical
posterior, after lifting \(q\) with its recognition evaluator and
retaining every evaluator sample. Then, for every \(t\in[\theta,2s]\),
\begin{align}
 \KL(q_t^E\Vert\mu_{\mathrm T,p(t)})&\le\kappa_E,
 \label{eq:phil-conditioned-full-entropy}\\
 \mathcal H_{\mathrm{opt}}(q_t^E)&\le\kappa_E,
 \label{eq:phil-conditioned-opt-entropy}\\
 \mathcal H_{\mathrm{phys}}(q_t^E)
 &\le\kappa_E+\log\frac{4+5\mathrm T}{4+\mathrm T}
 \le\kappa_E+\mathrm T.
 \label{eq:phil-conditioned-physical-entropy}
\end{align}
The two hidden-conditional bounds hold with any strictly positive
alternative source rows in their reference laws; the full-law bound
\cref{eq:phil-conditioned-full-entropy} refers to the native rows
\(p(t)\).

These are unscaled, dimensionless relative entropies. The corresponding
scaled optimizer free-energy loss is at most \(\mathrm T\kappa_E\);
the unscaled physical negative-ELBO loss is at most
\(\kappa_E+\mathrm T\), and its \(\mathrm T\)-scaled value is at most
\(\mathrm T\kappa_E+\mathrm T^2\).
Along the joint selection schedule with \(\eta,\delta\le\mathrm T\),
\[
 \rho=O(\mathrm T^{1/4}),
 \qquad
 \kappa_E=O(\mathrm T^{1/4}\log(1/\mathrm T)).
\]
The assumption \(\rho<1/2\) then holds for all sufficiently small
\(\mathrm T\), so these bounds are asymptotic.
\status{ESTABLISHED}

\noindent\emph{Proof.}
The learned-arm persistence bound
\cref{eq:phil-selected-coherence-persistence}, applied with
\(t_0=s\) and \(\mathsf h=1\), gives \(r:=\Pr(E^c)\le\rho\); its
proof uses only the learned-arm hypothesis. Thus \(E\) is a
\(\theta\)-measurable event of probability at least \(1-\rho\).
The law \(q_\theta^E\) is the actual conditional law and may have
retained--hidden correlations and a non-Gaussian hidden conditional.
Its marked pushforward is normalized by construction. At later times
no further event is conditioned upon. In particular, the deterministic
native row path is not recomputed from \(q_t^E\).

For a fixed \(t\in[\theta,2s]\), put
\(B=\mathbf1_E\) and \(R_t=(\mathbf X_t,S_t)\).
Compute the relative entropy of the joint law of \((B,R_t)\) against
\(\mathrm{Bernoulli}(1-r)\otimes\mu_{\mathrm T,p(t)}\) in two orders:
\begin{equation}
 \begin{aligned}
 &(1-r)\KL(q_t^E\Vert\mu_{\mathrm T,p(t)})
       +r\KL(q_t^{E^c}\Vert\mu_{\mathrm T,p(t)})\\
 &\qquad=\KL(q_t\Vert\mu_{\mathrm T,p(t)})
                       +\mathbb I(B;R_t)
 \le\eta+h_b(r).
 \end{aligned}
 \label{eq:phil-conditioned-event-information}
\end{equation}
When \(r=0\), use just the first term. The second term is nonnegative,
\(h_b\) is increasing on \([0,1/2]\), and \(1-r\ge1-\rho\).
This proves \cref{eq:phil-conditioned-full-entropy}.
It also establishes the needed absolute continuity and finiteness;
directly, \(q_t^E\le q_t/(1-r)\).

The channel's fibers can be checked explicitly. On a matching with
edges \(e,f\), let \(F_M\) have the orthonormal columns
\[
 (n_e^+-n_f^+)/\sqrt2,\qquad n_e^-,\qquad n_f^-.
\]
Write \(Z_M=\operatorname{im}F_M\), \(W_M=Z_M^\perp\), and
\(\mathbf x=y+F_Mz\).
The normal identities in \cref{eq:rg-formation-cycle-normals} give
\[
 \nu_{\mathrm T,Z}=\mathcal N(0,\mathrm TK_Z^{-1}),
 \qquad
 K_Z=\operatorname{diag}(8+\mathrm T,4+5\mathrm T,4+5\mathrm T)
\]
as the hidden optimizer conditional. Both the interaction Hessian
and the evaluator precision preserve \(W_M\) and \(Z_M\).
Consequently this conditional is independent of the retained
coordinate and of the prior rows. It is a reference conditional, not
an assertion that the actual conditioned recognition law has that
form. The same statement holds on the other matching by rotation;
off the two matchings the channel is the identity and has no
discarded conditional.

The target of \(C\) is a finite disjoint union of Euclidean strata,
so disintegration gives the exact entropy chain rule
\begin{equation}
 \KL(q_t^E\Vert\mu_{\mathrm T,p(t)})
 =\KL(C_\#q_t^E\Vert C_\#\mu_{\mathrm T,p(t)})
                         +\mathcal H_{\mathrm{opt}}(q_t^E).
 \label{eq:phil-conditioned-parent-chain}
\end{equation}
Dropping its retained term proves
\cref{eq:phil-conditioned-opt-entropy}. This argument uses the
complete actual law; it neither resets the hidden state nor assumes
Gaussianity conditional on a source history.

For the physical bound, work first on one matching fiber. Its physical
outer hidden Gaussian is
\[
 \Pi_{\mathrm T,Z}=\mathcal N(0,\mathrm T(K_Z^{\mathrm{phys}})^{-1}),
 \qquad
 K_Z^{\mathrm{phys}}
  =\operatorname{diag}(8+\mathrm T,4+\mathrm T,4+\mathrm T).
\]
All displayed matrices are positive definite. Write
\(v=1-\sigma_{\mathrm{ev}}^2=1/2\),
\(d_y=\mathfrak Dy\), and \(d_z=\mathfrak DF_Mz\).
Since \(R_{\mathrm{ev}}=\mathfrak D^\top\mathfrak D/v\)
preserves the split, \(d_y\cdot d_z=0\); its hidden eigenvalues
are \(0,4,4\).

Let \(U=(m,k)\) denote the full evaluator sample vectors.
At configuration \(\mathbf x\), the physical evaluator has
\(m\sim\mathcal N(\bar m,\sigma_{\mathrm{ev}}^2I)\) and independent
increment \(k-m\sim\mathcal N(0,vI)\). Its recognition counterpart,
\cref{eq:rg-formation-evaluator-lift}, has the same first law and
increment mean \(d_y+d_z\). At a fixed retained \(y\), form the
comparison probability law
\begin{equation}
 \overline P_y(dz,dU)
 =\Pi_{\mathrm T,Z}(dz)P_{\mathrm{ev}}(dU\mid y+F_Mz)
  \exp\left[
  \frac{d_y\cdot(k-m)-|d_y|^2/2}{v}\right].
 \label{eq:phil-conditioned-evaluator-tilt}
\end{equation}
Completing the Gaussian square in \(k-m\) proves that the multiplier
integrates to one for every \(z\). It changes the increment mean
to \(d_y\) and keeps \(\Pi_{\mathrm T,Z}\) as its \(z\) marginal.
The multiplier is retained-measurable in \((y,U)\), so
\(\overline P_y\) has the physical hidden conditional given these
retained variables. At a fixed \(z\), the recognition evaluator
differs by the additional increment mean \(d_z\), with relative
entropy \(|d_z|^2/(2v)\).

Compute the joint relative entropy first through \(z\), then through
\(U\), and drop the nonnegative evaluator-marginal term. For the
actual hidden law \(q(dz\mid M,y)\), this gives
\[
 \mathcal H_{\mathrm{phys}}(q\mid M,y)
 \le\KL(q(dz\mid M,y)\Vert\Pi_{\mathrm T,Z})
             +\E_q[|d_z|^2/(2v)\mid M,y].
\]
All terms are well-defined: the native process has finite second
moments, conditioning increases the second-moment bound by a factor
at most \((1-\rho)^{-1}\), and the Gaussian reference logarithms are quadratic.
The optimizer conditional is the exact tilt
\[
 \nu_{\mathrm T,Z}(dz)
 =\frac{\Pi_{\mathrm T,Z}(dz)e^{-|d_z|^2/(2v)}}{Z_{\mathrm{tilt}}},
 \qquad
 Z_{\mathrm{tilt}}=\frac{4+\mathrm T}{4+5\mathrm T}.
\]
For its normalizer, the common coordinate is unchanged and each of
the two discrepancy coordinates changes rate from \(4+\mathrm T\)
to \(4+5\mathrm T\); their two square-root Gaussian determinant
factors multiply to \(Z_{\mathrm{tilt}}\). Therefore
\[
 \KL(q(dz\mid M,y)\Vert\Pi_{\mathrm T,Z})
       +\E_q[|d_z|^2/(2v)\mid M,y]
 =\KL(q(dz\mid M,y)\Vert\nu_{\mathrm T,Z})
                         +\log\frac{4+5\mathrm T}{4+\mathrm T}.
\]
Integrate over the retained law. Matching strata have total
probability at most one, and all other strata have zero discarded
loss. Finally,
\[
 0\le\log\frac{4+5\mathrm T}{4+\mathrm T}
 =\log\left(1+\frac{4\mathrm T}{4+\mathrm T}\right)
 \le\frac{4\mathrm T}{4+\mathrm T}\le\mathrm T.
\]
This proves \cref{eq:phil-conditioned-physical-entropy}.
Only recognition has been conditioned on \(E\); the physical target
still uses the original observed records, with all evaluator samples
retained.

A source-prior factor \(\pi_S\) cancels in any conditional law on a
fiber because \(S\) is retained. The hidden optimizer and physical
reference conditionals are consequently the same for every strictly
positive choice of prior rows. The full-law entropy need not remain
small for different rows, because its retained term may change.
An imposed zero prior on an occupied stratum would remove reference
support and is excluded; ordinary finite-time learning from positive
rows preserves positivity.

The exact free-energy chain rules give the stated scaled and unscaled
losses. With \(\eta,\delta\le\mathrm T\) and
\(\mathsf a=\mathrm T^{1/4}\),
\[
 \rho\le3\mathrm T+8\sqrt{2/\pi}\mathrm T^{1/4}
                         +3408\mathrm T^{1/2}.
\]
Use \(h_b(\rho)\le\rho\log(e/\rho)\) and
\(1-\rho\ge1/2\) eventually to obtain the displayed order for
\(\kappa_E\). \(\square\)

\remarkheading{A high-probability prefix is needed}{rem:phil-conditioned-hidden-sign}
The high probability of the declared prefix has mathematical content.
Under the exact matching Gaussian \(\nu_{\mathrm T}^M\), a hidden
coordinate \(z_1\) is independent of the retained state and has zero
unconditioned hidden relative entropy. Conditioning on
\(\{z_1>0\}\), of probability \(1/2\), instead gives a density ratio
two on that half-space and hidden relative entropy \(\log2\) at every
temperature. Thus a probability merely bounded below does not suffice
for a vanishing conditional loss. Conditioning on the entire realized
fine history is different again: it fixes the hidden endpoint and
gives a Dirac conditional, singular relative to the nondegenerate
Gaussian fiber. The activation here uses the full joint law conditional
on \(E\), not a realized fine history.
\status{ESTABLISHED}

The proposition supplies conditional information at activation and for
the native future marginals under that same completed prefix. It does
not yet compare an independently evolving parent path. Its prefix is
the late restored coherence event, and carries no first-formation or
recursive-coalescence conclusion.



\subsection*{Inherited dynamics from the conditioned activation law}

\propositionheading{Dynamics of the conditioned parent}{prop:phil-actual-parent-dynamics}
Use the selected matching, channel, and conditioned activation of
\cref{prop:phil-actual-conditioned-parent}. More explicitly, let
\(0<\mathrm T\le1/100\), \(s=t_{\mathrm L}\), \(\theta=s+1\), and let \(\mathscr E\)
be the selected coherence event on \([s,\theta]\) at resolution
\(\mathsf a\) from \cref{prop:phil-selected-pair-persistence}.
Assume \(\Pr(\mathscr E^c)\le\rho<1\), and, throughout
\([s,\theta+\mathsf h]\subseteq[t_{\mathrm L},2t_{\mathrm L}]\),
\[
 \operatorname{TV}(q_t,\nu_{\mathrm T}^M)\le\delta,
 \qquad \max_i|p_i(t)-p_i^M|\le\delta.
\]
Write \(\varepsilon=\epsilon_{\mathrm L}\),
\(d=\delta+\varepsilon\mathsf h<1/4\), and use relative future time
\(u=t-\theta\). Initialize the parent with
\[
 q_0^C=C_\#\operatorname{Law}(\mathbf X_\theta,S_\theta\mid\mathscr E),
 \qquad p^C(0)=p(\theta).
\]
On the stratum \(\{S\}\times W_S\), its continuous generator is
\(-\mathsf A_S y\cdot\nabla_{W_S}+\mathrm T\Delta_{W_S}\).
Its rate-one proposal kernel is
\begin{equation}
 J_p((S,y),\cdot)=
 \int\kappa_S(d\mathbf x\mid y)
     \sum_{S'}b_p(S'\mid\mathbf x)
                   \delta_{(S',P_{W_{S'}}\mathbf x)},
 \label{eq:phil-actual-parent-kernel}
\end{equation}
where \(b_p(S'\mid\mathbf x)=\mu_{\mathrm T,p}(S'\mid\mathbf x)\) is the
product of receiver posterior rows at rows \(p\), and
\(\kappa_S\) is the optimizing hidden conditional: the three
independent Gaussian coordinates on a matching stratum, and a point
mass off the matchings. The parent rows evolve from its own law:
\begin{equation}
 \frac d{du}p_i^C(u)
 =\varepsilon\bigl(q_u^C(\sigma_i=+1)-p_i^C(u)\bigr).
 \label{eq:phil-actual-parent-learning}
\end{equation}
This coupled parent system has a unique global law with interior rows.

Let \(\mathsf P_{\theta,\mathsf h}^{\mathrm{fine},C}\) be the law of
\((C(\mathbf X_{\theta+u},S_{\theta+u}))_{0\le u\le\mathsf h}\)
conditioned only on \(\mathscr E\), with the fine process continuing
its original row path. Let \(\mathsf P_{\mathsf h}^C\) be the parent
path law from the displayed activation and its own learning equation.
Then
\begin{equation}
 \operatorname{TV}
 (\mathsf P_{\theta,\mathsf h}^{\mathrm{fine},C},
                                      \mathsf P_{\mathsf h}^C)
 \le2\rho+(\mathsf h+3)\delta+
                         \frac{4\mathsf h d}{1-4d}.
 \label{eq:phil-actual-parent-path-error}
\end{equation}
The comparison includes all retained coordinates and the assignment
history, hence also the loop history. Self-proposals remain part of
the rate-one kernel.

Along \cref{cor:phil-joint-selection}, choose
\(\delta\le\mathrm T\), \(\varepsilon\le\mathrm T^2\),
\(\mathsf a=\mathrm T^{1/4}\), and
\(\mathsf h=\mathrm T^{-1/4}\). One may take
\(\rho=3\delta+\mathsf b(\mathrm T,\mathsf a,1)\), where
\(\mathsf b\) is from \cref{prop:phil-selected-pair-persistence}.
The right side of \cref{eq:phil-actual-parent-path-error} is then
\(O(\mathrm T^{1/4})\). This is a conditional, finite-window
approximation after the selected coherence event; the activation law
is not replaced by a Gaussian law.
\status{ESTABLISHED}

\noindent\emph{Proof.}
For a matching \(M\), the decomposition
\(\R^{12}=W_M\oplus Z_M\) of
\cref{prop:rg-formation-endogenous-cycle} is preserved by
\(H_M\), \(\Omega\), and \(\mathsf A_M\).
The hidden drift matrix is
\(\operatorname{diag}(8+\mathrm T,4+5\mathrm T,4+5\mathrm T)\).
The prior weight \(\pi_S(p)\) cancels from the conditional law at
fixed \((S,y)\), so \(\kappa_S\) is independent of \(p\).
Off the matchings no coordinate is discarded. Thus the displayed
continuous generator is the projection of every held-sheet diffusion,
and \cref{eq:phil-actual-parent-kernel} is a probability kernel.
A self-proposal has \(S'=S\) and
\(P_{W_S}\mathbf x=y\), so it leaves the parent state unchanged.
At a different assignment it uses that assignment's projection,
including its twelve-dimensional stratum when appropriate.

For fixed rows, the pushed optimizing law \(C_\#\mu_{\mathrm T,p}\)
is invariant. Indeed each diffusion is reversible on its Gaussian
stratum, and
\(\mu_{\mathrm T,p}(d\mathbf x,S)b_p(S'\mid\mathbf x)\)
is symmetric in \(S,S'\). Pushing both endpoints through \(C\)
gives detailed balance for \(J_p\).
The loop covector annihilates each \(Z_M\), so this construction
retains the loop through the stratum changes as well.

We next check the existence assertion with evolving parent rows.
For any continuous interior row path, concatenate the linear
diffusions with the rate-one proposals, drawing \(\kappa_S\) at
each proposal. On every finite interval the clock has finitely many
rings and every inserted Gaussian is finite almost surely, so this
construction is nonexplosive. Its second moments remain finite:
projection cannot increase the norm, an inserted hidden draw adds
expected squared norm at most \(5\mathrm T/8\), and the diffusion
generator applied to \(|y|^2\) is at most
\(24\mathrm T-2\mathrm T|y|^2\). The full generator is consequently
bounded above by \(25\mathrm T-2\mathrm T|y|^2\).
The activation law has finite second moment because it conditions a
finite-moment fine law on an event of positive probability.

If two row paths have entries in \([\mathrm c_0,1-\mathrm c_0]\), couple their
parent processes by the same initial state, Brownian increments,
clocks, hidden Gaussian draws, and receiver uniforms until their
first different proposal. The derivative of a posterior receiver
probability with respect to its prior entry is at most
\(1/[4\mathrm c_0(1-\mathrm c_0)]\). Therefore
\[
 \operatorname{TV}(q_u^{C,p},q_u^{C,\widetilde p})
 \le\frac1{\mathrm c_0(1-\mathrm c_0)}
       \int_0^u\|p(v)-\widetilde p(v)\|_\infty dv.
\]
The variation-of-constants map for
\cref{eq:phil-actual-parent-learning} consequently has Lipschitz
constant at most \(\varepsilon u_*^2/[2\mathrm c_0(1-\mathrm c_0)]\) on an interval
of length \(u_*\). On a short enough interval it is a contraction
on continuous row paths in an interior cube. Its image stays in
that cube after further shortening the interval, since the row
velocity has absolute value at most \(\varepsilon\).
Moreover its solution satisfies
\[
 p_i^C(u)\ge p_i^C(0)e^{-\varepsilon u},\qquad
 1-p_i^C(u)\ge(1-p_i^C(0))e^{-\varepsilon u}.
\]
On any prescribed finite horizon these lower bounds give a common
interior cube and a positive contraction step length. Successive
steps give existence and uniqueness on that horizon, and then on
all finite horizons. The source marginal is continuous because its
jump rate is bounded, so the resulting rows are continuously
differentiable. This constructs learning from the parent's own law,
without imposing the fine row path after activation.

Let \(\overline\nu_{\mathrm T}^M\) be the retained Gaussian on
\(W_M\). Write \(\mathsf O^C_{\mathrm T,\mathsf h}\) for its
stationary held-\(M\) path law. The projection of
\(\mathsf O_{\mathrm T,\mathsf h}^M\) from
\cref{prop:phil-selected-pair-persistence} is exactly this law.
Only the learned-arm conclusions of that proposition are used, and
their proof does not use its frozen-arm hypothesis.
Conditioning any probability law on an event of probability
\(1-r\) changes its total variation by at most \(r\) after
any measurable projection. Thus the conditioned native retained
future obeys
\begin{equation}
 \operatorname{TV}
 (\mathsf P_{\theta,\mathsf h}^{\mathrm{fine},C},
                     \mathsf O^C_{\mathrm T,\mathsf h})
 \le\rho+(\mathsf h+2)\delta.
 \label{eq:phil-native-conditioned-held}
\end{equation}
This applies to a prefix event through \(\theta\). No condition is
placed on future coherence or future proposals. The same argument
at the activation time gives
\(\operatorname{TV}(q_0^C,
 \overline\nu_{\mathrm T}^M\otimes\delta_M)\le\rho+\delta\).

For the parent's own rows let \(\zeta_i^C\) be the probability in
row \(i\) assigned to the nonpartner of \(M\). Its equation and the
bound on any probability give
\[
 0\le\zeta_i^C(u)\le\zeta_i^C(0)+\varepsilon u
                  \le\delta+\varepsilon u\le d.
\]
The row-independent normalized sheet weights are at most the matching
weight one for \(0<\mathrm T\le1/100\) by
\cref{prop:phil-temperature-tracking}; hence, for any fixed row vector
with this bound,
\[
 \mu_{\mathrm T,p}(M)\ge\pi_M(p)
               \ge1-\sum_i\zeta_i\ge1-4d.
\]
Define the parent departure probability at a proposal by
\(k_p(y)=J_p((M,y),\{S'\ne M\})\).
Detailed balance, and the fact that the optimizing conditional
retained law on \(M\) is \(\overline\nu_{\mathrm T}^M\)
independently of the rows, give
\begin{equation}
 \int k_p(y)\overline\nu_{\mathrm T}^M(dy)
 \le\frac{1-\mu_{\mathrm T,p}(M)}{\mu_{\mathrm T,p}(M)}
 \le\frac{4d}{1-4d}.
 \label{eq:phil-parent-averaged-departure}
\end{equation}
The first inequality bounds stationary flux from \(M\) to its
complement by the complement's stationary mass.

Maximally couple the parent activation with the stationary retained
matching law and use common Brownian increments, clocks, and parent
proposal draws up to the first parent departure. Before departure,
the continuous paths agree because the matching drift splits, and
self-proposals do not move the retained coordinate. The row path
\(p^C(u)\) constructed above is deterministic, even though it was
computed from the parent's own law. Under the stationary comparison
path, its marginal at every time is
\(\overline\nu_{\mathrm T}^M\). Bound the first-departure
compensator with its stopped indicator by the compensator evaluated
on that whole stationary path. Applying
\cref{eq:phil-parent-averaged-departure} at each deterministic
\(p^C(u)\) gives
\[
 \Pr(\text{parent departure before }\mathsf h
                         \text{ while initially coupled})
 \le\frac{4\mathsf h d}{1-4d}.
\]
No Gaussian assertion is made about the parent's law conditioned on
survival or about the native hidden variables. We have proved
\[
 \operatorname{TV}(\mathsf P_{\mathsf h}^C,
                         \mathsf O^C_{\mathrm T,\mathsf h})
 \le\rho+\delta+\frac{4\mathsf h d}{1-4d}.
\]
Combining this with \cref{eq:phil-native-conditioned-held} proves
\cref{eq:phil-actual-parent-path-error}.

For the displayed joint choices,
\(d\le\mathrm T+\mathrm T^{7/4}\le2\mathrm T\),
\(\mathsf h\delta\le\mathrm T^{3/4}\), and
\(\rho=O(\mathrm T^{1/4})\).
The activation and following window fit inside the selection window
because \(1+\mathrm T^{-1/4}\le t_{\mathrm L}\) for sufficiently
small temperature under the mobility cap. The departure term is
\(O(\mathrm T^{3/4})\), giving the stated total error.
The Gaussian comparison laws served only to bound that error; the
parent's activation and subsequent learning remain the laws in the
statement. \hfill\(\square\)



